% Môn: Vật lý
% Lớp: 10
% Chương: Chương VI. Chuyển động tròn
% Bài: L10C6 Bài 32. Lực hướng tâm và gia tốc hướng tâm
% Số câu: 162
% App đọc file này trực tiếp — sửa rồi Commit trên GitHub, bấm Đồng bộ GitHub .tex

% L10C6 Bài 32. Lực hướng tâm và gia tốc hướng tâm

% ===== Câu 1 =====
% ID: AIQ_L10C6_FULL_0163
% Mức: NB
\dangbt{Nhận biết gia tốc hướng tâm và lực hướng tâm}
\begin{ex}
% ID: AIQ_L10C6_FULL_0163
% Mức: NB
Vật chuyển động tròn đều với tốc độ $2$ m/s trên bán kính $1$ m. Tính độ lớn gia tốc hướng tâm.
\choice
{\True $4\ \mathrm{m/s^2}$}
{$2\ \mathrm{m/s^2}$}
{Không đủ dữ kiện để có giá trị này}
{Không xác định}
\loigiai{
$a_{ht}=v^2/r=2^2/1$. Suy ra kết quả $4\ \mathrm{m/s^2}$.
}
\end{ex}

% ===== Câu 2 =====
% ID: AIQ_L10C6_FULL_0164
% Mức: NB
\begin{ex}
% ID: AIQ_L10C6_FULL_0164
% Mức: NB
Vật chuyển động tròn đều với tốc độ $3$ m/s trên bán kính $2$ m. Tính độ lớn gia tốc hướng tâm.
\choice
{$1,5\ \mathrm{m/s^2}$}
{\True $4,5\ \mathrm{m/s^2}$}
{$18\ \mathrm{m/s^2}$}
{$3\ \mathrm{m/s^2}$}
\loigiai{
$a_{ht}=v^2/r=3^2/2$. Suy ra kết quả $4,5\ \mathrm{m/s^2}$.
}
\end{ex}

% ===== Câu 3 =====
% ID: AIQ_L10C6_FULL_0165
% Mức: TH
\begin{ex}
% ID: AIQ_L10C6_FULL_0165
% Mức: TH
Vật chuyển động tròn đều với tốc độ $4$ m/s trên bán kính $3$ m. Tính độ lớn gia tốc hướng tâm.
\choice
{$1,33\ \mathrm{m/s^2}$}
{$48\ \mathrm{m/s^2}$}
{\True $5,33\ \mathrm{m/s^2}$}
{$2,67\ \mathrm{m/s^2}$}
\loigiai{
$a_{ht}=v^2/r=4^2/3$. Suy ra kết quả $5,33\ \mathrm{m/s^2}$.
}
\end{ex}

% ===== Câu 4 =====
% ID: AIQ_L10C6_FULL_0166
% Mức: TH
\begin{ex}
% ID: AIQ_L10C6_FULL_0166
% Mức: TH
Vật chuyển động tròn đều với tốc độ $5$ m/s trên bán kính $4$ m. Tính độ lớn gia tốc hướng tâm.
\choice
{$1,25\ \mathrm{m/s^2}$}
{$100\ \mathrm{m/s^2}$}
{$2,5\ \mathrm{m/s^2}$}
{\True $6,25\ \mathrm{m/s^2}$}
\loigiai{
$a_{ht}=v^2/r=5^2/4$. Suy ra kết quả $6,25\ \mathrm{m/s^2}$.
}
\end{ex}

% ===== Câu 5 =====
% ID: AIQ_L10C6_FULL_0167
% Mức: VD
\begin{ex}
% ID: AIQ_L10C6_FULL_0167
% Mức: VD
Vật chuyển động tròn đều với tốc độ $6$ m/s trên bán kính $1$ m. Tính độ lớn gia tốc hướng tâm.
\choice
{\True $36\ \mathrm{m/s^2}$}
{$6\ \mathrm{m/s^2}$}
{Không đủ dữ kiện để có giá trị này}
{$12\ \mathrm{m/s^2}$}
\loigiai{
$a_{ht}=v^2/r=6^2/1$. Suy ra kết quả $36\ \mathrm{m/s^2}$.
}
\end{ex}

% ===== Câu 6 =====
% ID: AIQ_L10C6_FULL_0168
% Mức: VD
\begin{ex}
% ID: AIQ_L10C6_FULL_0168
% Mức: VD
Vật chuyển động tròn đều với tốc độ $7$ m/s trên bán kính $2$ m. Tính độ lớn gia tốc hướng tâm.
\choice
{$3,5\ \mathrm{m/s^2}$}
{\True $24,5\ \mathrm{m/s^2}$}
{$98\ \mathrm{m/s^2}$}
{$7\ \mathrm{m/s^2}$}
\loigiai{
$a_{ht}=v^2/r=7^2/2$. Suy ra kết quả $24,5\ \mathrm{m/s^2}$.
}
\end{ex}

% ===== Câu 7 =====
% ID: AIQ_L10C6_FULL_0169
% Mức: VD
\begin{ex}
% ID: AIQ_L10C6_FULL_0169
% Mức: VD
Vật chuyển động tròn đều với tốc độ $8$ m/s trên bán kính $3$ m. Tính độ lớn gia tốc hướng tâm.
\choice
{$2,67\ \mathrm{m/s^2}$}
{$192\ \mathrm{m/s^2}$}
{\True $21,33\ \mathrm{m/s^2}$}
{$5,33\ \mathrm{m/s^2}$}
\loigiai{
$a_{ht}=v^2/r=8^2/3$. Suy ra kết quả $21,33\ \mathrm{m/s^2}$.
}
\end{ex}

% ===== Câu 8 =====
% ID: AIQ_L10C6_FULL_0170
% Mức: VDC
\begin{ex}
% ID: AIQ_L10C6_FULL_0170
% Mức: VDC
Vật chuyển động tròn đều với tốc độ $9$ m/s trên bán kính $4$ m. Tính độ lớn gia tốc hướng tâm.
\choice
{$2,25\ \mathrm{m/s^2}$}
{$324\ \mathrm{m/s^2}$}
{$4,5\ \mathrm{m/s^2}$}
{\True $20,25\ \mathrm{m/s^2}$}
\loigiai{
$a_{ht}=v^2/r=9^2/4$. Suy ra kết quả $20,25\ \mathrm{m/s^2}$.
}
\end{ex}

% ===== Câu 9 =====
% ID: AIQ_L10C6_FULL_0171
% Mức: VDC
\begin{ex}
% ID: AIQ_L10C6_FULL_0171
% Mức: VDC
Vật chuyển động tròn đều với tốc độ $10$ m/s trên bán kính $1$ m. Tính độ lớn gia tốc hướng tâm.
\choice
{\True $100\ \mathrm{m/s^2}$}
{$10\ \mathrm{m/s^2}$}
{Không đủ dữ kiện để có giá trị này}
{$20\ \mathrm{m/s^2}$}
\loigiai{
$a_{ht}=v^2/r=10^2/1$. Suy ra kết quả $100\ \mathrm{m/s^2}$.
}
\end{ex}

% ===== Câu 10 =====
% ID: AIQ_L10C6_FULL_0172
% Mức: VDC
\begin{ex}
% ID: AIQ_L10C6_FULL_0172
% Mức: VDC
Vật chuyển động tròn đều với tốc độ $11$ m/s trên bán kính $2$ m. Tính độ lớn gia tốc hướng tâm.
\choice
{$5,5\ \mathrm{m/s^2}$}
{\True $60,5\ \mathrm{m/s^2}$}
{$242\ \mathrm{m/s^2}$}
{$11\ \mathrm{m/s^2}$}
\loigiai{
$a_{ht}=v^2/r=11^2/2$. Suy ra kết quả $60,5\ \mathrm{m/s^2}$.
}
\end{ex}

% ===== Câu 11 =====
% ID: AIQ_L10C6_FULL_0173
% Mức: NB
\begin{ex}
% ID: AIQ_L10C6_FULL_0173
% Mức: NB
Xét các phát biểu thuộc dạng “Nhận biết gia tốc hướng tâm và lực hướng tâm” ở tình huống số 1.
\choiceTF
{\True Gia tốc hướng tâm hướng vào tâm quỹ đạo.}
{\True Gia tốc hướng tâm làm đổi hướng vận tốc.}
{Chuyển động tròn đều có gia tốc bằng 0.}
{\True Lực hướng tâm là hợp lực theo phương bán kính.}
\loigiai{
\begin{itemchoice}
\item (Đúng) Gia tốc hướng tâm hướng vào tâm quỹ đạo. (Vì): Phù hợp với định nghĩa hoặc hệ thức vật lí. b) Đúng: Phù hợp với định nghĩa hoặc hệ thức vật lí. c) Sai: Không phù hợp với định nghĩa hoặc hệ thức vật lí. d) Đúng: Phù hợp với định nghĩa hoặc hệ thức vật lí
\item (Đúng) Gia tốc hướng tâm làm đổi hướng vận tốc.
\item (Sai) Chuyển động tròn đều có gia tốc bằng 0.
\item (Đúng) Lực hướng tâm là hợp lực theo phương bán kính.
\end{itemchoice}
}
\end{ex}

% ===== Câu 12 =====
% ID: AIQ_L10C6_FULL_0174
% Mức: TH
\begin{ex}
% ID: AIQ_L10C6_FULL_0174
% Mức: TH
Xét các phát biểu thuộc dạng “Nhận biết gia tốc hướng tâm và lực hướng tâm” ở tình huống số 2.
\choiceTF
{\True Gia tốc hướng tâm hướng vào tâm quỹ đạo.}
{\True Gia tốc hướng tâm làm đổi hướng vận tốc.}
{Chuyển động tròn đều có gia tốc bằng 0.}
{\True Lực hướng tâm là hợp lực theo phương bán kính.}
\loigiai{
\begin{itemchoice}
\item (Đúng) Gia tốc hướng tâm hướng vào tâm quỹ đạo. (Vì): Phù hợp với định nghĩa hoặc hệ thức vật lí. b) Đúng: Phù hợp với định nghĩa hoặc hệ thức vật lí. c) Sai: Không phù hợp với định nghĩa hoặc hệ thức vật lí. d) Đúng: Phù hợp với định nghĩa hoặc hệ thức vật lí
\item (Đúng) Gia tốc hướng tâm làm đổi hướng vận tốc.
\item (Sai) Chuyển động tròn đều có gia tốc bằng 0.
\item (Đúng) Lực hướng tâm là hợp lực theo phương bán kính.
\end{itemchoice}
}
\end{ex}

% ===== Câu 13 =====
% ID: AIQ_L10C6_FULL_0175
% Mức: TH
\begin{ex}
% ID: AIQ_L10C6_FULL_0175
% Mức: TH
Xét các phát biểu thuộc dạng “Nhận biết gia tốc hướng tâm và lực hướng tâm” ở tình huống số 3.
\choiceTF
{\True Gia tốc hướng tâm hướng vào tâm quỹ đạo.}
{\True Gia tốc hướng tâm làm đổi hướng vận tốc.}
{Chuyển động tròn đều có gia tốc bằng 0.}
{\True Lực hướng tâm là hợp lực theo phương bán kính.}
\loigiai{
\begin{itemchoice}
\item (Đúng) Gia tốc hướng tâm hướng vào tâm quỹ đạo. (Vì): Phù hợp với định nghĩa hoặc hệ thức vật lí. b) Đúng: Phù hợp với định nghĩa hoặc hệ thức vật lí. c) Sai: Không phù hợp với định nghĩa hoặc hệ thức vật lí. d) Đúng: Phù hợp với định nghĩa hoặc hệ thức vật lí
\item (Đúng) Gia tốc hướng tâm làm đổi hướng vận tốc.
\item (Sai) Chuyển động tròn đều có gia tốc bằng 0.
\item (Đúng) Lực hướng tâm là hợp lực theo phương bán kính.
\end{itemchoice}
}
\end{ex}

% ===== Câu 14 =====
% ID: AIQ_L10C6_FULL_0176
% Mức: VD
\begin{ex}
% ID: AIQ_L10C6_FULL_0176
% Mức: VD
Xét các phát biểu thuộc dạng “Nhận biết gia tốc hướng tâm và lực hướng tâm” ở tình huống số 4.
\choiceTF
{\True Gia tốc hướng tâm hướng vào tâm quỹ đạo.}
{\True Gia tốc hướng tâm làm đổi hướng vận tốc.}
{Chuyển động tròn đều có gia tốc bằng 0.}
{\True Lực hướng tâm là hợp lực theo phương bán kính.}
\loigiai{
\begin{itemchoice}
\item (Đúng) Gia tốc hướng tâm hướng vào tâm quỹ đạo. (Vì): Phù hợp với định nghĩa hoặc hệ thức vật lí. b) Đúng: Phù hợp với định nghĩa hoặc hệ thức vật lí. c) Sai: Không phù hợp với định nghĩa hoặc hệ thức vật lí. d) Đúng: Phù hợp với định nghĩa hoặc hệ thức vật lí
\item (Đúng) Gia tốc hướng tâm làm đổi hướng vận tốc.
\item (Sai) Chuyển động tròn đều có gia tốc bằng 0.
\item (Đúng) Lực hướng tâm là hợp lực theo phương bán kính.
\end{itemchoice}
}
\end{ex}

% ===== Câu 15 =====
% ID: AIQ_L10C6_FULL_0177
% Mức: VDC
\begin{ex}
% ID: AIQ_L10C6_FULL_0177
% Mức: VDC
Xét các phát biểu thuộc dạng “Nhận biết gia tốc hướng tâm và lực hướng tâm” ở tình huống số 5.
\choiceTF
{\True Gia tốc hướng tâm hướng vào tâm quỹ đạo.}
{\True Gia tốc hướng tâm làm đổi hướng vận tốc.}
{Chuyển động tròn đều có gia tốc bằng 0.}
{\True Lực hướng tâm là hợp lực theo phương bán kính.}
\loigiai{
\begin{itemchoice}
\item (Đúng) Gia tốc hướng tâm hướng vào tâm quỹ đạo. (Vì): Phù hợp với định nghĩa hoặc hệ thức vật lí. b) Đúng: Phù hợp với định nghĩa hoặc hệ thức vật lí. c) Sai: Không phù hợp với định nghĩa hoặc hệ thức vật lí. d) Đúng: Phù hợp với định nghĩa hoặc hệ thức vật lí
\item (Đúng) Gia tốc hướng tâm làm đổi hướng vận tốc.
\item (Sai) Chuyển động tròn đều có gia tốc bằng 0.
\item (Đúng) Lực hướng tâm là hợp lực theo phương bán kính.
\end{itemchoice}
}
\end{ex}

% ===== Câu 16 =====
% ID: AIQ_L10C6_FULL_0178
% Mức: TH
\begin{ex}
% ID: AIQ_L10C6_FULL_0178
% Mức: TH
Vật chuyển động tròn đều với tốc độ $12$ m/s trên bán kính $3$ m. Tính độ lớn gia tốc hướng tâm.
\shortans{48}
\loigiai{
$a_{ht}=v^2/r=12^2/3$. Suy ra kết quả $48\ \mathrm{m/s^2}$. Đáp án trả lời ngắn không ghi đơn vị.
}
\end{ex}

% ===== Câu 17 =====
% ID: AIQ_L10C6_FULL_0179
% Mức: TH
\begin{ex}
% ID: AIQ_L10C6_FULL_0179
% Mức: TH
Vật chuyển động tròn đều với tốc độ $13$ m/s trên bán kính $4$ m. Tính độ lớn gia tốc hướng tâm.
\shortans{42,25}
\loigiai{
$a_{ht}=v^2/r=13^2/4$. Suy ra kết quả $42,25\ \mathrm{m/s^2}$. Đáp án trả lời ngắn không ghi đơn vị.
}
\end{ex}

% ===== Câu 18 =====
% ID: AIQ_L10C6_FULL_0180
% Mức: VD
\begin{ex}
% ID: AIQ_L10C6_FULL_0180
% Mức: VD
Vật chuyển động tròn đều với tốc độ $14$ m/s trên bán kính $1$ m. Tính độ lớn gia tốc hướng tâm.
\shortans{196}
\loigiai{
$a_{ht}=v^2/r=14^2/1$. Suy ra kết quả $196\ \mathrm{m/s^2}$. Đáp án trả lời ngắn không ghi đơn vị.
}
\end{ex}

% ===== Câu 19 =====
% ID: AIQ_L10C6_FULL_0181
% Mức: VD
\begin{ex}
% ID: AIQ_L10C6_FULL_0181
% Mức: VD
Vật chuyển động tròn đều với tốc độ $15$ m/s trên bán kính $2$ m. Tính độ lớn gia tốc hướng tâm.
\shortans{112,5}
\loigiai{
$a_{ht}=v^2/r=15^2/2$. Suy ra kết quả $112,5\ \mathrm{m/s^2}$. Đáp án trả lời ngắn không ghi đơn vị.
}
\end{ex}

% ===== Câu 20 =====
% ID: AIQ_L10C6_FULL_0182
% Mức: VDC
\begin{ex}
% ID: AIQ_L10C6_FULL_0182
% Mức: VDC
Vật chuyển động tròn đều với tốc độ $16$ m/s trên bán kính $3$ m. Tính độ lớn gia tốc hướng tâm.
\shortans{85,33}
\loigiai{
$a_{ht}=v^2/r=16^2/3$. Suy ra kết quả $85,33\ \mathrm{m/s^2}$. Đáp án trả lời ngắn không ghi đơn vị.
}
\end{ex}

% ===== Câu 21 =====
% ID: AIQ_L10C6_FULL_0183
% Mức: VDC
\begin{ex}
% ID: AIQ_L10C6_FULL_0183
% Mức: VDC
Vật chuyển động tròn đều với tốc độ $17$ m/s trên bán kính $4$ m. Tính độ lớn gia tốc hướng tâm.
\shortans{72,25}
\loigiai{
$a_{ht}=v^2/r=17^2/4$. Suy ra kết quả $72,25\ \mathrm{m/s^2}$. Đáp án trả lời ngắn không ghi đơn vị.
}
\end{ex}

% ===== Câu 22 =====
% ID: AIQ_L10C6_FULL_0184
% Mức: TH
\begin{bt}
% ID: AIQ_L10C6_FULL_0184
% Mức: TH
Trình bày cách giải: Vật chuyển động tròn đều với tốc độ $18$ m/s trên bán kính $1$ m. Tính độ lớn gia tốc hướng tâm.
\loigiai{
$a_{ht}=v^2/r=18^2/1$. Suy ra kết quả $324\ \mathrm{m/s^2}$.
}
\end{bt}

% ===== Câu 23 =====
% ID: AIQ_L10C6_FULL_0185
% Mức: TH
\begin{bt}
% ID: AIQ_L10C6_FULL_0185
% Mức: TH
Trình bày cách giải: Vật chuyển động tròn đều với tốc độ $19$ m/s trên bán kính $2$ m. Tính độ lớn gia tốc hướng tâm.
\loigiai{
$a_{ht}=v^2/r=19^2/2$. Suy ra kết quả $180,5\ \mathrm{m/s^2}$.
}
\end{bt}

% ===== Câu 24 =====
% ID: AIQ_L10C6_FULL_0186
% Mức: VD
\begin{bt}
% ID: AIQ_L10C6_FULL_0186
% Mức: VD
Trình bày cách giải: Vật chuyển động tròn đều với tốc độ $20$ m/s trên bán kính $3$ m. Tính độ lớn gia tốc hướng tâm.
\loigiai{
$a_{ht}=v^2/r=20^2/3$. Suy ra kết quả $133,33\ \mathrm{m/s^2}$.
}
\end{bt}

% ===== Câu 25 =====
% ID: AIQ_L10C6_FULL_0187
% Mức: VD
\begin{bt}
% ID: AIQ_L10C6_FULL_0187
% Mức: VD
Trình bày cách giải: Vật chuyển động tròn đều với tốc độ $21$ m/s trên bán kính $4$ m. Tính độ lớn gia tốc hướng tâm.
\loigiai{
$a_{ht}=v^2/r=21^2/4$. Suy ra kết quả $110,25\ \mathrm{m/s^2}$.
}
\end{bt}

% ===== Câu 26 =====
% ID: AIQ_L10C6_FULL_0188
% Mức: VDC
\begin{bt}
% ID: AIQ_L10C6_FULL_0188
% Mức: VDC
Trình bày cách giải: Vật chuyển động tròn đều với tốc độ $22$ m/s trên bán kính $1$ m. Tính độ lớn gia tốc hướng tâm.
\loigiai{
$a_{ht}=v^2/r=22^2/1$. Suy ra kết quả $484\ \mathrm{m/s^2}$.
}
\end{bt}

% ===== Câu 27 =====
% ID: AIQ_L10C6_FULL_0189
% Mức: VDC
\begin{bt}
% ID: AIQ_L10C6_FULL_0189
% Mức: VDC
Trình bày cách giải: Vật chuyển động tròn đều với tốc độ $23$ m/s trên bán kính $2$ m. Tính độ lớn gia tốc hướng tâm.
\loigiai{
$a_{ht}=v^2/r=23^2/2$. Suy ra kết quả $264,5\ \mathrm{m/s^2}$.
}
\end{bt}

% ===== Câu 28 =====
% ID: AIQ_L10C6_FULL_0190
% Mức: NB
\dangbt{Tính gia tốc hướng tâm}
\begin{ex}
% ID: AIQ_L10C6_FULL_0190
% Mức: NB
Vật quay đều với tốc độ góc $2$ rad/s, bán kính $0,5$ m. Tính gia tốc hướng tâm.
\choice
{\True $2\ \mathrm{m/s^2}$}
{$1\ \mathrm{m/s^2}$}
{$8\ \mathrm{m/s^2}$}
{Không đủ dữ kiện để có giá trị này}
\loigiai{
$a_{ht}=\omega^2r=2^2\cdot0,5$. Suy ra kết quả $2\ \mathrm{m/s^2}$.
}
\end{ex}

% ===== Câu 29 =====
% ID: AIQ_L10C6_FULL_0191
% Mức: NB
\begin{ex}
% ID: AIQ_L10C6_FULL_0191
% Mức: NB
Vật quay đều với tốc độ góc $3$ rad/s, bán kính $0,6$ m. Tính gia tốc hướng tâm.
\choice
{$1,8\ \mathrm{m/s^2}$}
{\True $5,4\ \mathrm{m/s^2}$}
{$15\ \mathrm{m/s^2}$}
{$3,6\ \mathrm{m/s^2}$}
\loigiai{
$a_{ht}=\omega^2r=3^2\cdot0,6$. Suy ra kết quả $5,4\ \mathrm{m/s^2}$.
}
\end{ex}

% ===== Câu 30 =====
% ID: AIQ_L10C6_FULL_0192
% Mức: TH
\begin{ex}
% ID: AIQ_L10C6_FULL_0192
% Mức: TH
Vật quay đều với tốc độ góc $4$ rad/s, bán kính $0,7$ m. Tính gia tốc hướng tâm.
\choice
{$2,8\ \mathrm{m/s^2}$}
{$22,86\ \mathrm{m/s^2}$}
{\True $11,2\ \mathrm{m/s^2}$}
{$5,6\ \mathrm{m/s^2}$}
\loigiai{
$a_{ht}=\omega^2r=4^2\cdot0,7$. Suy ra kết quả $11,2\ \mathrm{m/s^2}$.
}
\end{ex}

% ===== Câu 31 =====
% ID: AIQ_L10C6_FULL_0193
% Mức: TH
\begin{ex}
% ID: AIQ_L10C6_FULL_0193
% Mức: TH
Vật quay đều với tốc độ góc $5$ rad/s, bán kính $0,8$ m. Tính gia tốc hướng tâm.
\choice
{$4\ \mathrm{m/s^2}$}
{$31,25\ \mathrm{m/s^2}$}
{$8\ \mathrm{m/s^2}$}
{\True $20\ \mathrm{m/s^2}$}
\loigiai{
$a_{ht}=\omega^2r=5^2\cdot0,8$. Suy ra kết quả $20\ \mathrm{m/s^2}$.
}
\end{ex}

% ===== Câu 32 =====
% ID: AIQ_L10C6_FULL_0194
% Mức: VD
\begin{ex}
% ID: AIQ_L10C6_FULL_0194
% Mức: VD
Vật quay đều với tốc độ góc $6$ rad/s, bán kính $0,5$ m. Tính gia tốc hướng tâm.
\choice
{\True $18\ \mathrm{m/s^2}$}
{$3\ \mathrm{m/s^2}$}
{$72\ \mathrm{m/s^2}$}
{$6\ \mathrm{m/s^2}$}
\loigiai{
$a_{ht}=\omega^2r=6^2\cdot0,5$. Suy ra kết quả $18\ \mathrm{m/s^2}$.
}
\end{ex}

% ===== Câu 33 =====
% ID: AIQ_L10C6_FULL_0195
% Mức: VD
\begin{ex}
% ID: AIQ_L10C6_FULL_0195
% Mức: VD
Vật quay đều với tốc độ góc $2$ rad/s, bán kính $0,6$ m. Tính gia tốc hướng tâm.
\choice
{$1,2\ \mathrm{m/s^2}$}
{\True $2,4\ \mathrm{m/s^2}$}
{$6,67\ \mathrm{m/s^2}$}
{Không đủ dữ kiện để có giá trị này}
\loigiai{
$a_{ht}=\omega^2r=2^2\cdot0,6$. Suy ra kết quả $2,4\ \mathrm{m/s^2}$.
}
\end{ex}

% ===== Câu 34 =====
% ID: AIQ_L10C6_FULL_0196
% Mức: VD
\begin{ex}
% ID: AIQ_L10C6_FULL_0196
% Mức: VD
Vật quay đều với tốc độ góc $3$ rad/s, bán kính $0,7$ m. Tính gia tốc hướng tâm.
\choice
{$2,1\ \mathrm{m/s^2}$}
{$12,86\ \mathrm{m/s^2}$}
{\True $6,3\ \mathrm{m/s^2}$}
{$4,2\ \mathrm{m/s^2}$}
\loigiai{
$a_{ht}=\omega^2r=3^2\cdot0,7$. Suy ra kết quả $6,3\ \mathrm{m/s^2}$.
}
\end{ex}

% ===== Câu 35 =====
% ID: AIQ_L10C6_FULL_0197
% Mức: VDC
\begin{ex}
% ID: AIQ_L10C6_FULL_0197
% Mức: VDC
Vật quay đều với tốc độ góc $4$ rad/s, bán kính $0,8$ m. Tính gia tốc hướng tâm.
\choice
{$3,2\ \mathrm{m/s^2}$}
{$20\ \mathrm{m/s^2}$}
{$6,4\ \mathrm{m/s^2}$}
{\True $12,8\ \mathrm{m/s^2}$}
\loigiai{
$a_{ht}=\omega^2r=4^2\cdot0,8$. Suy ra kết quả $12,8\ \mathrm{m/s^2}$.
}
\end{ex}

% ===== Câu 36 =====
% ID: AIQ_L10C6_FULL_0198
% Mức: VDC
\begin{ex}
% ID: AIQ_L10C6_FULL_0198
% Mức: VDC
Vật quay đều với tốc độ góc $5$ rad/s, bán kính $0,5$ m. Tính gia tốc hướng tâm.
\choice
{\True $12,5\ \mathrm{m/s^2}$}
{$2,5\ \mathrm{m/s^2}$}
{$50\ \mathrm{m/s^2}$}
{$5\ \mathrm{m/s^2}$}
\loigiai{
$a_{ht}=\omega^2r=5^2\cdot0,5$. Suy ra kết quả $12,5\ \mathrm{m/s^2}$.
}
\end{ex}

% ===== Câu 37 =====
% ID: AIQ_L10C6_FULL_0199
% Mức: VDC
\begin{ex}
% ID: AIQ_L10C6_FULL_0199
% Mức: VDC
Vật quay đều với tốc độ góc $6$ rad/s, bán kính $0,6$ m. Tính gia tốc hướng tâm.
\choice
{$3,6\ \mathrm{m/s^2}$}
{\True $21,6\ \mathrm{m/s^2}$}
{$60\ \mathrm{m/s^2}$}
{$7,2\ \mathrm{m/s^2}$}
\loigiai{
$a_{ht}=\omega^2r=6^2\cdot0,6$. Suy ra kết quả $21,6\ \mathrm{m/s^2}$.
}
\end{ex}

% ===== Câu 38 =====
% ID: AIQ_L10C6_FULL_0200
% Mức: NB
\begin{ex}
% ID: AIQ_L10C6_FULL_0200
% Mức: NB
Xét các phát biểu thuộc dạng “Tính gia tốc hướng tâm” ở tình huống số 1.
\choiceTF
{\True $a_{ht}=v^2/r$.}
{\True $a_{ht}=\omega^2r$.}
{Đơn vị gia tốc là N.}
{\True Tăng tốc độ làm gia tốc hướng tâm tăng theo bình phương.}
\loigiai{
\begin{itemchoice}
\item (Đúng) $a_{ht}=v^2/r$. (Vì): Phù hợp với định nghĩa hoặc hệ thức vật lí. b) Đúng: Phù hợp với định nghĩa hoặc hệ thức vật lí. c) Sai: Không phù hợp với định nghĩa hoặc hệ thức vật lí. d) Đúng: Phù hợp với định nghĩa hoặc hệ thức vật lí
\item (Đúng) $a_{ht}=\omega^2r$.
\item (Sai) Đơn vị gia tốc là N.
\item (Đúng) Tăng tốc độ làm gia tốc hướng tâm tăng theo bình phương.
\end{itemchoice}
}
\end{ex}

% ===== Câu 39 =====
% ID: AIQ_L10C6_FULL_0201
% Mức: TH
\begin{ex}
% ID: AIQ_L10C6_FULL_0201
% Mức: TH
Xét các phát biểu thuộc dạng “Tính gia tốc hướng tâm” ở tình huống số 2.
\choiceTF
{\True $a_{ht}=v^2/r$.}
{\True $a_{ht}=\omega^2r$.}
{Đơn vị gia tốc là N.}
{\True Tăng tốc độ làm gia tốc hướng tâm tăng theo bình phương.}
\loigiai{
\begin{itemchoice}
\item (Đúng) $a_{ht}=v^2/r$. (Vì): Phù hợp với định nghĩa hoặc hệ thức vật lí. b) Đúng: Phù hợp với định nghĩa hoặc hệ thức vật lí. c) Sai: Không phù hợp với định nghĩa hoặc hệ thức vật lí. d) Đúng: Phù hợp với định nghĩa hoặc hệ thức vật lí
\item (Đúng) $a_{ht}=\omega^2r$.
\item (Sai) Đơn vị gia tốc là N.
\item (Đúng) Tăng tốc độ làm gia tốc hướng tâm tăng theo bình phương.
\end{itemchoice}
}
\end{ex}

% ===== Câu 40 =====
% ID: AIQ_L10C6_FULL_0202
% Mức: TH
\begin{ex}
% ID: AIQ_L10C6_FULL_0202
% Mức: TH
Xét các phát biểu thuộc dạng “Tính gia tốc hướng tâm” ở tình huống số 3.
\choiceTF
{\True $a_{ht}=v^2/r$.}
{\True $a_{ht}=\omega^2r$.}
{Đơn vị gia tốc là N.}
{\True Tăng tốc độ làm gia tốc hướng tâm tăng theo bình phương.}
\loigiai{
\begin{itemchoice}
\item (Đúng) $a_{ht}=v^2/r$. (Vì): Phù hợp với định nghĩa hoặc hệ thức vật lí. b) Đúng: Phù hợp với định nghĩa hoặc hệ thức vật lí. c) Sai: Không phù hợp với định nghĩa hoặc hệ thức vật lí. d) Đúng: Phù hợp với định nghĩa hoặc hệ thức vật lí
\item (Đúng) $a_{ht}=\omega^2r$.
\item (Sai) Đơn vị gia tốc là N.
\item (Đúng) Tăng tốc độ làm gia tốc hướng tâm tăng theo bình phương.
\end{itemchoice}
}
\end{ex}

% ===== Câu 41 =====
% ID: AIQ_L10C6_FULL_0203
% Mức: VD
\begin{ex}
% ID: AIQ_L10C6_FULL_0203
% Mức: VD
Xét các phát biểu thuộc dạng “Tính gia tốc hướng tâm” ở tình huống số 4.
\choiceTF
{\True $a_{ht}=v^2/r$.}
{\True $a_{ht}=\omega^2r$.}
{Đơn vị gia tốc là N.}
{\True Tăng tốc độ làm gia tốc hướng tâm tăng theo bình phương.}
\loigiai{
\begin{itemchoice}
\item (Đúng) $a_{ht}=v^2/r$. (Vì): Phù hợp với định nghĩa hoặc hệ thức vật lí. b) Đúng: Phù hợp với định nghĩa hoặc hệ thức vật lí. c) Sai: Không phù hợp với định nghĩa hoặc hệ thức vật lí. d) Đúng: Phù hợp với định nghĩa hoặc hệ thức vật lí
\item (Đúng) $a_{ht}=\omega^2r$.
\item (Sai) Đơn vị gia tốc là N.
\item (Đúng) Tăng tốc độ làm gia tốc hướng tâm tăng theo bình phương.
\end{itemchoice}
}
\end{ex}

% ===== Câu 42 =====
% ID: AIQ_L10C6_FULL_0204
% Mức: VDC
\begin{ex}
% ID: AIQ_L10C6_FULL_0204
% Mức: VDC
Xét các phát biểu thuộc dạng “Tính gia tốc hướng tâm” ở tình huống số 5.
\choiceTF
{\True $a_{ht}=v^2/r$.}
{\True $a_{ht}=\omega^2r$.}
{Đơn vị gia tốc là N.}
{\True Tăng tốc độ làm gia tốc hướng tâm tăng theo bình phương.}
\loigiai{
\begin{itemchoice}
\item (Đúng) $a_{ht}=v^2/r$. (Vì): Phù hợp với định nghĩa hoặc hệ thức vật lí. b) Đúng: Phù hợp với định nghĩa hoặc hệ thức vật lí. c) Sai: Không phù hợp với định nghĩa hoặc hệ thức vật lí. d) Đúng: Phù hợp với định nghĩa hoặc hệ thức vật lí
\item (Đúng) $a_{ht}=\omega^2r$.
\item (Sai) Đơn vị gia tốc là N.
\item (Đúng) Tăng tốc độ làm gia tốc hướng tâm tăng theo bình phương.
\end{itemchoice}
}
\end{ex}

% ===== Câu 43 =====
% ID: AIQ_L10C6_FULL_0205
% Mức: TH
\begin{ex}
% ID: AIQ_L10C6_FULL_0205
% Mức: TH
Vật quay đều với tốc độ góc $2$ rad/s, bán kính $0,7$ m. Tính gia tốc hướng tâm.
\shortans{2,8}
\loigiai{
$a_{ht}=\omega^2r=2^2\cdot0,7$. Suy ra kết quả $2,8\ \mathrm{m/s^2}$. Đáp án trả lời ngắn không ghi đơn vị.
}
\end{ex}

% ===== Câu 44 =====
% ID: AIQ_L10C6_FULL_0206
% Mức: TH
\begin{ex}
% ID: AIQ_L10C6_FULL_0206
% Mức: TH
Vật quay đều với tốc độ góc $3$ rad/s, bán kính $0,8$ m. Tính gia tốc hướng tâm.
\shortans{7,2}
\loigiai{
$a_{ht}=\omega^2r=3^2\cdot0,8$. Suy ra kết quả $7,2\ \mathrm{m/s^2}$. Đáp án trả lời ngắn không ghi đơn vị.
}
\end{ex}

% ===== Câu 45 =====
% ID: AIQ_L10C6_FULL_0207
% Mức: VD
\begin{ex}
% ID: AIQ_L10C6_FULL_0207
% Mức: VD
Vật quay đều với tốc độ góc $4$ rad/s, bán kính $0,5$ m. Tính gia tốc hướng tâm.
\shortans{8}
\loigiai{
$a_{ht}=\omega^2r=4^2\cdot0,5$. Suy ra kết quả $8\ \mathrm{m/s^2}$. Đáp án trả lời ngắn không ghi đơn vị.
}
\end{ex}

% ===== Câu 46 =====
% ID: AIQ_L10C6_FULL_0208
% Mức: VD
\begin{ex}
% ID: AIQ_L10C6_FULL_0208
% Mức: VD
Vật quay đều với tốc độ góc $5$ rad/s, bán kính $0,6$ m. Tính gia tốc hướng tâm.
\shortans{15}
\loigiai{
$a_{ht}=\omega^2r=5^2\cdot0,6$. Suy ra kết quả $15\ \mathrm{m/s^2}$. Đáp án trả lời ngắn không ghi đơn vị.
}
\end{ex}

% ===== Câu 47 =====
% ID: AIQ_L10C6_FULL_0209
% Mức: VDC
\begin{ex}
% ID: AIQ_L10C6_FULL_0209
% Mức: VDC
Vật quay đều với tốc độ góc $6$ rad/s, bán kính $0,7$ m. Tính gia tốc hướng tâm.
\shortans{25,2}
\loigiai{
$a_{ht}=\omega^2r=6^2\cdot0,7$. Suy ra kết quả $25,2\ \mathrm{m/s^2}$. Đáp án trả lời ngắn không ghi đơn vị.
}
\end{ex}

% ===== Câu 48 =====
% ID: AIQ_L10C6_FULL_0210
% Mức: VDC
\begin{ex}
% ID: AIQ_L10C6_FULL_0210
% Mức: VDC
Vật quay đều với tốc độ góc $2$ rad/s, bán kính $0,8$ m. Tính gia tốc hướng tâm.
\shortans{3,2}
\loigiai{
$a_{ht}=\omega^2r=2^2\cdot0,8$. Suy ra kết quả $3,2\ \mathrm{m/s^2}$. Đáp án trả lời ngắn không ghi đơn vị.
}
\end{ex}

% ===== Câu 49 =====
% ID: AIQ_L10C6_FULL_0211
% Mức: TH
\begin{bt}
% ID: AIQ_L10C6_FULL_0211
% Mức: TH
Trình bày cách giải: Vật quay đều với tốc độ góc $3$ rad/s, bán kính $0,5$ m. Tính gia tốc hướng tâm.
\loigiai{
$a_{ht}=\omega^2r=3^2\cdot0,5$. Suy ra kết quả $4,5\ \mathrm{m/s^2}$.
}
\end{bt}

% ===== Câu 50 =====
% ID: AIQ_L10C6_FULL_0212
% Mức: TH
\begin{bt}
% ID: AIQ_L10C6_FULL_0212
% Mức: TH
Trình bày cách giải: Vật quay đều với tốc độ góc $4$ rad/s, bán kính $0,6$ m. Tính gia tốc hướng tâm.
\loigiai{
$a_{ht}=\omega^2r=4^2\cdot0,6$. Suy ra kết quả $9,6\ \mathrm{m/s^2}$.
}
\end{bt}

% ===== Câu 51 =====
% ID: AIQ_L10C6_FULL_0213
% Mức: VD
\begin{bt}
% ID: AIQ_L10C6_FULL_0213
% Mức: VD
Trình bày cách giải: Vật quay đều với tốc độ góc $5$ rad/s, bán kính $0,7$ m. Tính gia tốc hướng tâm.
\loigiai{
$a_{ht}=\omega^2r=5^2\cdot0,7$. Suy ra kết quả $17,5\ \mathrm{m/s^2}$.
}
\end{bt}

% ===== Câu 52 =====
% ID: AIQ_L10C6_FULL_0214
% Mức: VD
\begin{bt}
% ID: AIQ_L10C6_FULL_0214
% Mức: VD
Trình bày cách giải: Vật quay đều với tốc độ góc $6$ rad/s, bán kính $0,8$ m. Tính gia tốc hướng tâm.
\loigiai{
$a_{ht}=\omega^2r=6^2\cdot0,8$. Suy ra kết quả $28,8\ \mathrm{m/s^2}$.
}
\end{bt}

% ===== Câu 53 =====
% ID: AIQ_L10C6_FULL_0215
% Mức: VDC
\begin{bt}
% ID: AIQ_L10C6_FULL_0215
% Mức: VDC
Trình bày cách giải: Vật quay đều với tốc độ góc $2$ rad/s, bán kính $0,5$ m. Tính gia tốc hướng tâm.
\loigiai{
$a_{ht}=\omega^2r=2^2\cdot0,5$. Suy ra kết quả $2\ \mathrm{m/s^2}$.
}
\end{bt}

% ===== Câu 54 =====
% ID: AIQ_L10C6_FULL_0216
% Mức: VDC
\begin{bt}
% ID: AIQ_L10C6_FULL_0216
% Mức: VDC
Trình bày cách giải: Vật quay đều với tốc độ góc $3$ rad/s, bán kính $0,6$ m. Tính gia tốc hướng tâm.
\loigiai{
$a_{ht}=\omega^2r=3^2\cdot0,6$. Suy ra kết quả $5,4\ \mathrm{m/s^2}$.
}
\end{bt}

% ===== Câu 55 =====
% ID: AIQ_L10C6_FULL_0217
% Mức: NB
\dangbt{Tính lực hướng tâm trong chuyển động tròn đều}
\begin{ex}
% ID: AIQ_L10C6_FULL_0217
% Mức: NB
Vật khối lượng $0,2$ kg chuyển động tròn đều tốc độ $3$ m/s, bán kính $1$ m. Tính lực hướng tâm.
\choice
{\True $1,8\ \mathrm{N}$}
{$0,6\ \mathrm{N}$}
{Không đủ dữ kiện để có giá trị này}
{$9\ \mathrm{N}$}
\loigiai{
$F_{ht}=mv^2/r=0,2\cdot3^2/1$. Suy ra kết quả $1,8\ \mathrm{N}$.
}
\end{ex}

% ===== Câu 56 =====
% ID: AIQ_L10C6_FULL_0218
% Mức: NB
\begin{ex}
% ID: AIQ_L10C6_FULL_0218
% Mức: NB
Vật khối lượng $0,3$ kg chuyển động tròn đều tốc độ $4$ m/s, bán kính $2$ m. Tính lực hướng tâm.
\choice
{$0,6\ \mathrm{N}$}
{\True $2,4\ \mathrm{N}$}
{$9,6\ \mathrm{N}$}
{$8\ \mathrm{N}$}
\loigiai{
$F_{ht}=mv^2/r=0,3\cdot4^2/2$. Suy ra kết quả $2,4\ \mathrm{N}$.
}
\end{ex}

% ===== Câu 57 =====
% ID: AIQ_L10C6_FULL_0219
% Mức: TH
\begin{ex}
% ID: AIQ_L10C6_FULL_0219
% Mức: TH
Vật khối lượng $0,4$ kg chuyển động tròn đều tốc độ $5$ m/s, bán kính $3$ m. Tính lực hướng tâm.
\choice
{$0,67\ \mathrm{N}$}
{$30\ \mathrm{N}$}
{\True $3,33\ \mathrm{N}$}
{$8,33\ \mathrm{N}$}
\loigiai{
$F_{ht}=mv^2/r=0,4\cdot5^2/3$. Suy ra kết quả $3,33\ \mathrm{N}$.
}
\end{ex}

% ===== Câu 58 =====
% ID: AIQ_L10C6_FULL_0220
% Mức: TH
\begin{ex}
% ID: AIQ_L10C6_FULL_0220
% Mức: TH
Vật khối lượng $0,5$ kg chuyển động tròn đều tốc độ $6$ m/s, bán kính $4$ m. Tính lực hướng tâm.
\choice
{$0,75\ \mathrm{N}$}
{$72\ \mathrm{N}$}
{$9\ \mathrm{N}$}
{\True $4,5\ \mathrm{N}$}
\loigiai{
$F_{ht}=mv^2/r=0,5\cdot6^2/4$. Suy ra kết quả $4,5\ \mathrm{N}$.
}
\end{ex}

% ===== Câu 59 =====
% ID: AIQ_L10C6_FULL_0221
% Mức: VD
\begin{ex}
% ID: AIQ_L10C6_FULL_0221
% Mức: VD
Vật khối lượng $0,6$ kg chuyển động tròn đều tốc độ $7$ m/s, bán kính $1$ m. Tính lực hướng tâm.
\choice
{\True $29,4\ \mathrm{N}$}
{$4,2\ \mathrm{N}$}
{Không đủ dữ kiện để có giá trị này}
{$49\ \mathrm{N}$}
\loigiai{
$F_{ht}=mv^2/r=0,6\cdot7^2/1$. Suy ra kết quả $29,4\ \mathrm{N}$.
}
\end{ex}

% ===== Câu 60 =====
% ID: AIQ_L10C6_FULL_0222
% Mức: VD
\begin{ex}
% ID: AIQ_L10C6_FULL_0222
% Mức: VD
Vật khối lượng $0,2$ kg chuyển động tròn đều tốc độ $8$ m/s, bán kính $2$ m. Tính lực hướng tâm.
\choice
{$0,8\ \mathrm{N}$}
{\True $6,4\ \mathrm{N}$}
{$25,6\ \mathrm{N}$}
{$32\ \mathrm{N}$}
\loigiai{
$F_{ht}=mv^2/r=0,2\cdot8^2/2$. Suy ra kết quả $6,4\ \mathrm{N}$.
}
\end{ex}

% ===== Câu 61 =====
% ID: AIQ_L10C6_FULL_0223
% Mức: VD
\begin{ex}
% ID: AIQ_L10C6_FULL_0223
% Mức: VD
Vật khối lượng $0,3$ kg chuyển động tròn đều tốc độ $9$ m/s, bán kính $3$ m. Tính lực hướng tâm.
\choice
{$0,9\ \mathrm{N}$}
{$72,9\ \mathrm{N}$}
{\True $8,1\ \mathrm{N}$}
{$27\ \mathrm{N}$}
\loigiai{
$F_{ht}=mv^2/r=0,3\cdot9^2/3$. Suy ra kết quả $8,1\ \mathrm{N}$.
}
\end{ex}

% ===== Câu 62 =====
% ID: AIQ_L10C6_FULL_0224
% Mức: VDC
\begin{ex}
% ID: AIQ_L10C6_FULL_0224
% Mức: VDC
Vật khối lượng $0,4$ kg chuyển động tròn đều tốc độ $10$ m/s, bán kính $4$ m. Tính lực hướng tâm.
\choice
{$1\ \mathrm{N}$}
{$160\ \mathrm{N}$}
{$25\ \mathrm{N}$}
{\True $10\ \mathrm{N}$}
\loigiai{
$F_{ht}=mv^2/r=0,4\cdot10^2/4$. Suy ra kết quả $10\ \mathrm{N}$.
}
\end{ex}

% ===== Câu 63 =====
% ID: AIQ_L10C6_FULL_0225
% Mức: VDC
\begin{ex}
% ID: AIQ_L10C6_FULL_0225
% Mức: VDC
Vật khối lượng $0,5$ kg chuyển động tròn đều tốc độ $11$ m/s, bán kính $1$ m. Tính lực hướng tâm.
\choice
{\True $60,5\ \mathrm{N}$}
{$5,5\ \mathrm{N}$}
{Không đủ dữ kiện để có giá trị này}
{$121\ \mathrm{N}$}
\loigiai{
$F_{ht}=mv^2/r=0,5\cdot11^2/1$. Suy ra kết quả $60,5\ \mathrm{N}$.
}
\end{ex}

% ===== Câu 64 =====
% ID: AIQ_L10C6_FULL_0226
% Mức: VDC
\begin{ex}
% ID: AIQ_L10C6_FULL_0226
% Mức: VDC
Vật khối lượng $0,6$ kg chuyển động tròn đều tốc độ $12$ m/s, bán kính $2$ m. Tính lực hướng tâm.
\choice
{$3,6\ \mathrm{N}$}
{\True $43,2\ \mathrm{N}$}
{$172,8\ \mathrm{N}$}
{$72\ \mathrm{N}$}
\loigiai{
$F_{ht}=mv^2/r=0,6\cdot12^2/2$. Suy ra kết quả $43,2\ \mathrm{N}$.
}
\end{ex}

% ===== Câu 65 =====
% ID: AIQ_L10C6_FULL_0227
% Mức: NB
\begin{ex}
% ID: AIQ_L10C6_FULL_0227
% Mức: NB
Xét các phát biểu thuộc dạng “Tính lực hướng tâm trong chuyển động tròn đều” ở tình huống số 1.
\choiceTF
{\True $F_{ht}=mv^2/r$.}
{Lực hướng tâm là một lực mới ngoài các lực đã có.}
{\True Hợp lực hướng tâm có thể do lực căng tạo ra.}
{\True Đơn vị lực hướng tâm là N.}
\loigiai{
\begin{itemchoice}
\item (Đúng) $F_{ht}=mv^2/r$. (Vì): Phù hợp với định nghĩa hoặc hệ thức vật lí. b) Sai: Không phù hợp với định nghĩa hoặc hệ thức vật lí. c) Đúng: Phù hợp với định nghĩa hoặc hệ thức vật lí. d) Đúng: Phù hợp với định nghĩa hoặc hệ thức vật lí
\item (Sai) Lực hướng tâm là một lực mới ngoài các lực đã có.
\item (Đúng) Hợp lực hướng tâm có thể do lực căng tạo ra.
\item (Đúng) Đơn vị lực hướng tâm là N.
\end{itemchoice}
}
\end{ex}

% ===== Câu 66 =====
% ID: AIQ_L10C6_FULL_0228
% Mức: TH
\begin{ex}
% ID: AIQ_L10C6_FULL_0228
% Mức: TH
Xét các phát biểu thuộc dạng “Tính lực hướng tâm trong chuyển động tròn đều” ở tình huống số 2.
\choiceTF
{\True $F_{ht}=mv^2/r$.}
{Lực hướng tâm là một lực mới ngoài các lực đã có.}
{\True Hợp lực hướng tâm có thể do lực căng tạo ra.}
{\True Đơn vị lực hướng tâm là N.}
\loigiai{
\begin{itemchoice}
\item (Đúng) $F_{ht}=mv^2/r$. (Vì): Phù hợp với định nghĩa hoặc hệ thức vật lí. b) Sai: Không phù hợp với định nghĩa hoặc hệ thức vật lí. c) Đúng: Phù hợp với định nghĩa hoặc hệ thức vật lí. d) Đúng: Phù hợp với định nghĩa hoặc hệ thức vật lí
\item (Sai) Lực hướng tâm là một lực mới ngoài các lực đã có.
\item (Đúng) Hợp lực hướng tâm có thể do lực căng tạo ra.
\item (Đúng) Đơn vị lực hướng tâm là N.
\end{itemchoice}
}
\end{ex}

% ===== Câu 67 =====
% ID: AIQ_L10C6_FULL_0229
% Mức: TH
\begin{ex}
% ID: AIQ_L10C6_FULL_0229
% Mức: TH
Xét các phát biểu thuộc dạng “Tính lực hướng tâm trong chuyển động tròn đều” ở tình huống số 3.
\choiceTF
{\True $F_{ht}=mv^2/r$.}
{Lực hướng tâm là một lực mới ngoài các lực đã có.}
{\True Hợp lực hướng tâm có thể do lực căng tạo ra.}
{\True Đơn vị lực hướng tâm là N.}
\loigiai{
\begin{itemchoice}
\item (Đúng) $F_{ht}=mv^2/r$. (Vì): Phù hợp với định nghĩa hoặc hệ thức vật lí. b) Sai: Không phù hợp với định nghĩa hoặc hệ thức vật lí. c) Đúng: Phù hợp với định nghĩa hoặc hệ thức vật lí. d) Đúng: Phù hợp với định nghĩa hoặc hệ thức vật lí
\item (Sai) Lực hướng tâm là một lực mới ngoài các lực đã có.
\item (Đúng) Hợp lực hướng tâm có thể do lực căng tạo ra.
\item (Đúng) Đơn vị lực hướng tâm là N.
\end{itemchoice}
}
\end{ex}

% ===== Câu 68 =====
% ID: AIQ_L10C6_FULL_0230
% Mức: VD
\begin{ex}
% ID: AIQ_L10C6_FULL_0230
% Mức: VD
Xét các phát biểu thuộc dạng “Tính lực hướng tâm trong chuyển động tròn đều” ở tình huống số 4.
\choiceTF
{\True $F_{ht}=mv^2/r$.}
{Lực hướng tâm là một lực mới ngoài các lực đã có.}
{\True Hợp lực hướng tâm có thể do lực căng tạo ra.}
{\True Đơn vị lực hướng tâm là N.}
\loigiai{
\begin{itemchoice}
\item (Đúng) $F_{ht}=mv^2/r$. (Vì): Phù hợp với định nghĩa hoặc hệ thức vật lí. b) Sai: Không phù hợp với định nghĩa hoặc hệ thức vật lí. c) Đúng: Phù hợp với định nghĩa hoặc hệ thức vật lí. d) Đúng: Phù hợp với định nghĩa hoặc hệ thức vật lí
\item (Sai) Lực hướng tâm là một lực mới ngoài các lực đã có.
\item (Đúng) Hợp lực hướng tâm có thể do lực căng tạo ra.
\item (Đúng) Đơn vị lực hướng tâm là N.
\end{itemchoice}
}
\end{ex}

% ===== Câu 69 =====
% ID: AIQ_L10C6_FULL_0231
% Mức: VDC
\begin{ex}
% ID: AIQ_L10C6_FULL_0231
% Mức: VDC
Xét các phát biểu thuộc dạng “Tính lực hướng tâm trong chuyển động tròn đều” ở tình huống số 5.
\choiceTF
{\True $F_{ht}=mv^2/r$.}
{Lực hướng tâm là một lực mới ngoài các lực đã có.}
{\True Hợp lực hướng tâm có thể do lực căng tạo ra.}
{\True Đơn vị lực hướng tâm là N.}
\loigiai{
\begin{itemchoice}
\item (Đúng) $F_{ht}=mv^2/r$. (Vì): Phù hợp với định nghĩa hoặc hệ thức vật lí. b) Sai: Không phù hợp với định nghĩa hoặc hệ thức vật lí. c) Đúng: Phù hợp với định nghĩa hoặc hệ thức vật lí. d) Đúng: Phù hợp với định nghĩa hoặc hệ thức vật lí
\item (Sai) Lực hướng tâm là một lực mới ngoài các lực đã có.
\item (Đúng) Hợp lực hướng tâm có thể do lực căng tạo ra.
\item (Đúng) Đơn vị lực hướng tâm là N.
\end{itemchoice}
}
\end{ex}

% ===== Câu 70 =====
% ID: AIQ_L10C6_FULL_0232
% Mức: TH
\begin{ex}
% ID: AIQ_L10C6_FULL_0232
% Mức: TH
Vật khối lượng $0,2$ kg chuyển động tròn đều tốc độ $13$ m/s, bán kính $3$ m. Tính lực hướng tâm.
\shortans{11,27}
\loigiai{
$F_{ht}=mv^2/r=0,2\cdot13^2/3$. Suy ra kết quả $11,27\ \mathrm{N}$. Đáp án trả lời ngắn không ghi đơn vị.
}
\end{ex}

% ===== Câu 71 =====
% ID: AIQ_L10C6_FULL_0233
% Mức: TH
\begin{ex}
% ID: AIQ_L10C6_FULL_0233
% Mức: TH
Vật khối lượng $0,3$ kg chuyển động tròn đều tốc độ $14$ m/s, bán kính $4$ m. Tính lực hướng tâm.
\shortans{14,7}
\loigiai{
$F_{ht}=mv^2/r=0,3\cdot14^2/4$. Suy ra kết quả $14,7\ \mathrm{N}$. Đáp án trả lời ngắn không ghi đơn vị.
}
\end{ex}

% ===== Câu 72 =====
% ID: AIQ_L10C6_FULL_0234
% Mức: VD
\begin{ex}
% ID: AIQ_L10C6_FULL_0234
% Mức: VD
Vật khối lượng $0,4$ kg chuyển động tròn đều tốc độ $15$ m/s, bán kính $1$ m. Tính lực hướng tâm.
\shortans{90}
\loigiai{
$F_{ht}=mv^2/r=0,4\cdot15^2/1$. Suy ra kết quả $90\ \mathrm{N}$. Đáp án trả lời ngắn không ghi đơn vị.
}
\end{ex}

% ===== Câu 73 =====
% ID: AIQ_L10C6_FULL_0235
% Mức: VD
\begin{ex}
% ID: AIQ_L10C6_FULL_0235
% Mức: VD
Vật khối lượng $0,5$ kg chuyển động tròn đều tốc độ $16$ m/s, bán kính $2$ m. Tính lực hướng tâm.
\shortans{64}
\loigiai{
$F_{ht}=mv^2/r=0,5\cdot16^2/2$. Suy ra kết quả $64\ \mathrm{N}$. Đáp án trả lời ngắn không ghi đơn vị.
}
\end{ex}

% ===== Câu 74 =====
% ID: AIQ_L10C6_FULL_0236
% Mức: VDC
\begin{ex}
% ID: AIQ_L10C6_FULL_0236
% Mức: VDC
Vật khối lượng $0,6$ kg chuyển động tròn đều tốc độ $17$ m/s, bán kính $3$ m. Tính lực hướng tâm.
\shortans{57,8}
\loigiai{
$F_{ht}=mv^2/r=0,6\cdot17^2/3$. Suy ra kết quả $57,8\ \mathrm{N}$. Đáp án trả lời ngắn không ghi đơn vị.
}
\end{ex}

% ===== Câu 75 =====
% ID: AIQ_L10C6_FULL_0237
% Mức: VDC
\begin{ex}
% ID: AIQ_L10C6_FULL_0237
% Mức: VDC
Vật khối lượng $0,2$ kg chuyển động tròn đều tốc độ $18$ m/s, bán kính $4$ m. Tính lực hướng tâm.
\shortans{16,2}
\loigiai{
$F_{ht}=mv^2/r=0,2\cdot18^2/4$. Suy ra kết quả $16,2\ \mathrm{N}$. Đáp án trả lời ngắn không ghi đơn vị.
}
\end{ex}

% ===== Câu 76 =====
% ID: AIQ_L10C6_FULL_0238
% Mức: TH
\begin{bt}
% ID: AIQ_L10C6_FULL_0238
% Mức: TH
Trình bày cách giải: Vật khối lượng $0,3$ kg chuyển động tròn đều tốc độ $19$ m/s, bán kính $1$ m. Tính lực hướng tâm.
\loigiai{
$F_{ht}=mv^2/r=0,3\cdot19^2/1$. Suy ra kết quả $108,3\ \mathrm{N}$.
}
\end{bt}

% ===== Câu 77 =====
% ID: AIQ_L10C6_FULL_0239
% Mức: TH
\begin{bt}
% ID: AIQ_L10C6_FULL_0239
% Mức: TH
Trình bày cách giải: Vật khối lượng $0,4$ kg chuyển động tròn đều tốc độ $20$ m/s, bán kính $2$ m. Tính lực hướng tâm.
\loigiai{
$F_{ht}=mv^2/r=0,4\cdot20^2/2$. Suy ra kết quả $80\ \mathrm{N}$.
}
\end{bt}

% ===== Câu 78 =====
% ID: AIQ_L10C6_FULL_0240
% Mức: VD
\begin{bt}
% ID: AIQ_L10C6_FULL_0240
% Mức: VD
Trình bày cách giải: Vật khối lượng $0,5$ kg chuyển động tròn đều tốc độ $21$ m/s, bán kính $3$ m. Tính lực hướng tâm.
\loigiai{
$F_{ht}=mv^2/r=0,5\cdot21^2/3$. Suy ra kết quả $73,5\ \mathrm{N}$.
}
\end{bt}

% ===== Câu 79 =====
% ID: AIQ_L10C6_FULL_0241
% Mức: VD
\begin{bt}
% ID: AIQ_L10C6_FULL_0241
% Mức: VD
Trình bày cách giải: Vật khối lượng $0,6$ kg chuyển động tròn đều tốc độ $22$ m/s, bán kính $4$ m. Tính lực hướng tâm.
\loigiai{
$F_{ht}=mv^2/r=0,6\cdot22^2/4$. Suy ra kết quả $72,6\ \mathrm{N}$.
}
\end{bt}

% ===== Câu 80 =====
% ID: AIQ_L10C6_FULL_0242
% Mức: VDC
\begin{bt}
% ID: AIQ_L10C6_FULL_0242
% Mức: VDC
Trình bày cách giải: Vật khối lượng $0,2$ kg chuyển động tròn đều tốc độ $23$ m/s, bán kính $1$ m. Tính lực hướng tâm.
\loigiai{
$F_{ht}=mv^2/r=0,2\cdot23^2/1$. Suy ra kết quả $105,8\ \mathrm{N}$.
}
\end{bt}

% ===== Câu 81 =====
% ID: AIQ_L10C6_FULL_0243
% Mức: VDC
\begin{bt}
% ID: AIQ_L10C6_FULL_0243
% Mức: VDC
Trình bày cách giải: Vật khối lượng $0,3$ kg chuyển động tròn đều tốc độ $24$ m/s, bán kính $2$ m. Tính lực hướng tâm.
\loigiai{
$F_{ht}=mv^2/r=0,3\cdot24^2/2$. Suy ra kết quả $86,4\ \mathrm{N}$.
}
\end{bt}

% ===== Câu 82 =====
% ID: AIQ_L10C6_FULL_0244
% Mức: NB
\dangbt{Chuyển động của vật trên mặt phẳng ngang và đường cong}
\begin{ex}
% ID: AIQ_L10C6_FULL_0244
% Mức: NB
Ô tô qua khúc cua phẳng bán kính $20$ m với tốc độ $6$ m/s. Hệ số ma sát nghỉ tối thiểu, lấy $g=10$ m/s², bằng bao nhiêu?
\choice
{\True $0,18\ \mathrm{}$}
{$0,03\ \mathrm{}$}
{$1,8\ \mathrm{}$}
{$18\ \mathrm{}$}
\loigiai{
$\mu mg=mv^2/r\Rightarrow\mu=v^2/(gr)=6^2/(10\cdot20)$. Suy ra kết quả $0,18\ \mathrm{}$.
}
\end{ex}

% ===== Câu 83 =====
% ID: AIQ_L10C6_FULL_0245
% Mức: NB
\begin{ex}
% ID: AIQ_L10C6_FULL_0245
% Mức: NB
Ô tô qua khúc cua phẳng bán kính $22$ m với tốc độ $7$ m/s. Hệ số ma sát nghỉ tối thiểu, lấy $g=10$ m/s², bằng bao nhiêu?
\choice
{$0,03\ \mathrm{}$}
{\True $0,22\ \mathrm{}$}
{$2,23\ \mathrm{}$}
{$22,27\ \mathrm{}$}
\loigiai{
$\mu mg=mv^2/r\Rightarrow\mu=v^2/(gr)=7^2/(10\cdot22)$. Suy ra kết quả $0,22\ \mathrm{}$.
}
\end{ex}

% ===== Câu 84 =====
% ID: AIQ_L10C6_FULL_0246
% Mức: TH
\begin{ex}
% ID: AIQ_L10C6_FULL_0246
% Mức: TH
Ô tô qua khúc cua phẳng bán kính $24$ m với tốc độ $8$ m/s. Hệ số ma sát nghỉ tối thiểu, lấy $g=10$ m/s², bằng bao nhiêu?
\choice
{$0,03\ \mathrm{}$}
{$2,67\ \mathrm{}$}
{\True $0,27\ \mathrm{}$}
{$26,67\ \mathrm{}$}
\loigiai{
$\mu mg=mv^2/r\Rightarrow\mu=v^2/(gr)=8^2/(10\cdot24)$. Suy ra kết quả $0,27\ \mathrm{}$.
}
\end{ex}

% ===== Câu 85 =====
% ID: AIQ_L10C6_FULL_0247
% Mức: TH
\begin{ex}
% ID: AIQ_L10C6_FULL_0247
% Mức: TH
Ô tô qua khúc cua phẳng bán kính $26$ m với tốc độ $9$ m/s. Hệ số ma sát nghỉ tối thiểu, lấy $g=10$ m/s², bằng bao nhiêu?
\choice
{$0,03\ \mathrm{}$}
{$3,12\ \mathrm{}$}
{$31,15\ \mathrm{}$}
{\True $0,31\ \mathrm{}$}
\loigiai{
$\mu mg=mv^2/r\Rightarrow\mu=v^2/(gr)=9^2/(10\cdot26)$. Suy ra kết quả $0,31\ \mathrm{}$.
}
\end{ex}

% ===== Câu 86 =====
% ID: AIQ_L10C6_FULL_0248
% Mức: VD
\begin{ex}
% ID: AIQ_L10C6_FULL_0248
% Mức: VD
Ô tô qua khúc cua phẳng bán kính $28$ m với tốc độ $10$ m/s. Hệ số ma sát nghỉ tối thiểu, lấy $g=10$ m/s², bằng bao nhiêu?
\choice
{\True $0,36\ \mathrm{}$}
{$0,04\ \mathrm{}$}
{$3,57\ \mathrm{}$}
{$35,71\ \mathrm{}$}
\loigiai{
$\mu mg=mv^2/r\Rightarrow\mu=v^2/(gr)=10^2/(10\cdot28)$. Suy ra kết quả $0,36\ \mathrm{}$.
}
\end{ex}

% ===== Câu 87 =====
% ID: AIQ_L10C6_FULL_0249
% Mức: VD
\begin{ex}
% ID: AIQ_L10C6_FULL_0249
% Mức: VD
Ô tô qua khúc cua phẳng bán kính $30$ m với tốc độ $11$ m/s. Hệ số ma sát nghỉ tối thiểu, lấy $g=10$ m/s², bằng bao nhiêu?
\choice
{$0,04\ \mathrm{}$}
{\True $0,4\ \mathrm{}$}
{$4,03\ \mathrm{}$}
{$40,33\ \mathrm{}$}
\loigiai{
$\mu mg=mv^2/r\Rightarrow\mu=v^2/(gr)=11^2/(10\cdot30)$. Suy ra kết quả $0,4\ \mathrm{}$.
}
\end{ex}

% ===== Câu 88 =====
% ID: AIQ_L10C6_FULL_0250
% Mức: VD
\begin{ex}
% ID: AIQ_L10C6_FULL_0250
% Mức: VD
Ô tô qua khúc cua phẳng bán kính $32$ m với tốc độ $12$ m/s. Hệ số ma sát nghỉ tối thiểu, lấy $g=10$ m/s², bằng bao nhiêu?
\choice
{$0,04\ \mathrm{}$}
{$4,5\ \mathrm{}$}
{\True $0,45\ \mathrm{}$}
{$45\ \mathrm{}$}
\loigiai{
$\mu mg=mv^2/r\Rightarrow\mu=v^2/(gr)=12^2/(10\cdot32)$. Suy ra kết quả $0,45\ \mathrm{}$.
}
\end{ex}

% ===== Câu 89 =====
% ID: AIQ_L10C6_FULL_0251
% Mức: VDC
\begin{ex}
% ID: AIQ_L10C6_FULL_0251
% Mức: VDC
Ô tô qua khúc cua phẳng bán kính $34$ m với tốc độ $13$ m/s. Hệ số ma sát nghỉ tối thiểu, lấy $g=10$ m/s², bằng bao nhiêu?
\choice
{$0,04\ \mathrm{}$}
{$4,97\ \mathrm{}$}
{$49,71\ \mathrm{}$}
{\True $0,5\ \mathrm{}$}
\loigiai{
$\mu mg=mv^2/r\Rightarrow\mu=v^2/(gr)=13^2/(10\cdot34)$. Suy ra kết quả $0,5\ \mathrm{}$.
}
\end{ex}

% ===== Câu 90 =====
% ID: AIQ_L10C6_FULL_0252
% Mức: VDC
\begin{ex}
% ID: AIQ_L10C6_FULL_0252
% Mức: VDC
Ô tô qua khúc cua phẳng bán kính $36$ m với tốc độ $14$ m/s. Hệ số ma sát nghỉ tối thiểu, lấy $g=10$ m/s², bằng bao nhiêu?
\choice
{\True $0,54\ \mathrm{}$}
{$0,04\ \mathrm{}$}
{$5,44\ \mathrm{}$}
{$54,44\ \mathrm{}$}
\loigiai{
$\mu mg=mv^2/r\Rightarrow\mu=v^2/(gr)=14^2/(10\cdot36)$. Suy ra kết quả $0,54\ \mathrm{}$.
}
\end{ex}

% ===== Câu 91 =====
% ID: AIQ_L10C6_FULL_0253
% Mức: VDC
\begin{ex}
% ID: AIQ_L10C6_FULL_0253
% Mức: VDC
Ô tô qua khúc cua phẳng bán kính $38$ m với tốc độ $15$ m/s. Hệ số ma sát nghỉ tối thiểu, lấy $g=10$ m/s², bằng bao nhiêu?
\choice
{$0,04\ \mathrm{}$}
{\True $0,59\ \mathrm{}$}
{$5,92\ \mathrm{}$}
{$59,21\ \mathrm{}$}
\loigiai{
$\mu mg=mv^2/r\Rightarrow\mu=v^2/(gr)=15^2/(10\cdot38)$. Suy ra kết quả $0,59\ \mathrm{}$.
}
\end{ex}

% ===== Câu 92 =====
% ID: AIQ_L10C6_FULL_0254
% Mức: NB
\begin{ex}
% ID: AIQ_L10C6_FULL_0254
% Mức: NB
Xét các phát biểu thuộc dạng “Chuyển động của vật trên mặt phẳng ngang và đường cong” ở tình huống số 1.
\choiceTF
{\True Ma sát có thể đóng vai trò lực hướng tâm khi xe vào cua phẳng.}
{\True Tốc độ càng lớn càng dễ trượt.}
{\True Bán kính cua lớn làm yêu cầu lực hướng tâm giảm.}
{Xe vào cua không cần lực hướng tâm.}
\loigiai{
\begin{itemchoice}
\item (Đúng) Ma sát có thể đóng vai trò lực hướng tâm khi xe vào cua phẳng. (Vì): Phù hợp với định nghĩa hoặc hệ thức vật lí. b) Đúng: Phù hợp với định nghĩa hoặc hệ thức vật lí. c) Đúng: Phù hợp với định nghĩa hoặc hệ thức vật lí. d) Sai: Không phù hợp với định nghĩa hoặc hệ thức vật lí
\item (Đúng) Tốc độ càng lớn càng dễ trượt.
\item (Đúng) Bán kính cua lớn làm yêu cầu lực hướng tâm giảm.
\item (Sai) Xe vào cua không cần lực hướng tâm.
\end{itemchoice}
}
\end{ex}

% ===== Câu 93 =====
% ID: AIQ_L10C6_FULL_0255
% Mức: TH
\begin{ex}
% ID: AIQ_L10C6_FULL_0255
% Mức: TH
Xét các phát biểu thuộc dạng “Chuyển động của vật trên mặt phẳng ngang và đường cong” ở tình huống số 2.
\choiceTF
{\True Ma sát có thể đóng vai trò lực hướng tâm khi xe vào cua phẳng.}
{\True Tốc độ càng lớn càng dễ trượt.}
{\True Bán kính cua lớn làm yêu cầu lực hướng tâm giảm.}
{Xe vào cua không cần lực hướng tâm.}
\loigiai{
\begin{itemchoice}
\item (Đúng) Ma sát có thể đóng vai trò lực hướng tâm khi xe vào cua phẳng. (Vì): Phù hợp với định nghĩa hoặc hệ thức vật lí. b) Đúng: Phù hợp với định nghĩa hoặc hệ thức vật lí. c) Đúng: Phù hợp với định nghĩa hoặc hệ thức vật lí. d) Sai: Không phù hợp với định nghĩa hoặc hệ thức vật lí
\item (Đúng) Tốc độ càng lớn càng dễ trượt.
\item (Đúng) Bán kính cua lớn làm yêu cầu lực hướng tâm giảm.
\item (Sai) Xe vào cua không cần lực hướng tâm.
\end{itemchoice}
}
\end{ex}

% ===== Câu 94 =====
% ID: AIQ_L10C6_FULL_0256
% Mức: TH
\begin{ex}
% ID: AIQ_L10C6_FULL_0256
% Mức: TH
Xét các phát biểu thuộc dạng “Chuyển động của vật trên mặt phẳng ngang và đường cong” ở tình huống số 3.
\choiceTF
{\True Ma sát có thể đóng vai trò lực hướng tâm khi xe vào cua phẳng.}
{\True Tốc độ càng lớn càng dễ trượt.}
{\True Bán kính cua lớn làm yêu cầu lực hướng tâm giảm.}
{Xe vào cua không cần lực hướng tâm.}
\loigiai{
\begin{itemchoice}
\item (Đúng) Ma sát có thể đóng vai trò lực hướng tâm khi xe vào cua phẳng. (Vì): Phù hợp với định nghĩa hoặc hệ thức vật lí. b) Đúng: Phù hợp với định nghĩa hoặc hệ thức vật lí. c) Đúng: Phù hợp với định nghĩa hoặc hệ thức vật lí. d) Sai: Không phù hợp với định nghĩa hoặc hệ thức vật lí
\item (Đúng) Tốc độ càng lớn càng dễ trượt.
\item (Đúng) Bán kính cua lớn làm yêu cầu lực hướng tâm giảm.
\item (Sai) Xe vào cua không cần lực hướng tâm.
\end{itemchoice}
}
\end{ex}

% ===== Câu 95 =====
% ID: AIQ_L10C6_FULL_0257
% Mức: VD
\begin{ex}
% ID: AIQ_L10C6_FULL_0257
% Mức: VD
Xét các phát biểu thuộc dạng “Chuyển động của vật trên mặt phẳng ngang và đường cong” ở tình huống số 4.
\choiceTF
{\True Ma sát có thể đóng vai trò lực hướng tâm khi xe vào cua phẳng.}
{\True Tốc độ càng lớn càng dễ trượt.}
{\True Bán kính cua lớn làm yêu cầu lực hướng tâm giảm.}
{Xe vào cua không cần lực hướng tâm.}
\loigiai{
\begin{itemchoice}
\item (Đúng) Ma sát có thể đóng vai trò lực hướng tâm khi xe vào cua phẳng. (Vì): Phù hợp với định nghĩa hoặc hệ thức vật lí. b) Đúng: Phù hợp với định nghĩa hoặc hệ thức vật lí. c) Đúng: Phù hợp với định nghĩa hoặc hệ thức vật lí. d) Sai: Không phù hợp với định nghĩa hoặc hệ thức vật lí
\item (Đúng) Tốc độ càng lớn càng dễ trượt.
\item (Đúng) Bán kính cua lớn làm yêu cầu lực hướng tâm giảm.
\item (Sai) Xe vào cua không cần lực hướng tâm.
\end{itemchoice}
}
\end{ex}

% ===== Câu 96 =====
% ID: AIQ_L10C6_FULL_0258
% Mức: VDC
\begin{ex}
% ID: AIQ_L10C6_FULL_0258
% Mức: VDC
Xét các phát biểu thuộc dạng “Chuyển động của vật trên mặt phẳng ngang và đường cong” ở tình huống số 5.
\choiceTF
{\True Ma sát có thể đóng vai trò lực hướng tâm khi xe vào cua phẳng.}
{\True Tốc độ càng lớn càng dễ trượt.}
{\True Bán kính cua lớn làm yêu cầu lực hướng tâm giảm.}
{Xe vào cua không cần lực hướng tâm.}
\loigiai{
\begin{itemchoice}
\item (Đúng) Ma sát có thể đóng vai trò lực hướng tâm khi xe vào cua phẳng. (Vì): Phù hợp với định nghĩa hoặc hệ thức vật lí. b) Đúng: Phù hợp với định nghĩa hoặc hệ thức vật lí. c) Đúng: Phù hợp với định nghĩa hoặc hệ thức vật lí. d) Sai: Không phù hợp với định nghĩa hoặc hệ thức vật lí
\item (Đúng) Tốc độ càng lớn càng dễ trượt.
\item (Đúng) Bán kính cua lớn làm yêu cầu lực hướng tâm giảm.
\item (Sai) Xe vào cua không cần lực hướng tâm.
\end{itemchoice}
}
\end{ex}

% ===== Câu 97 =====
% ID: AIQ_L10C6_FULL_0259
% Mức: TH
\begin{ex}
% ID: AIQ_L10C6_FULL_0259
% Mức: TH
Ô tô qua khúc cua phẳng bán kính $40$ m với tốc độ $16$ m/s. Hệ số ma sát nghỉ tối thiểu, lấy $g=10$ m/s², bằng bao nhiêu?
\shortans{0,64}
\loigiai{
$\mu mg=mv^2/r\Rightarrow\mu=v^2/(gr)=16^2/(10\cdot40)$. Suy ra kết quả $0,64\ \mathrm{}$. Đáp án trả lời ngắn không ghi đơn vị.
}
\end{ex}

% ===== Câu 98 =====
% ID: AIQ_L10C6_FULL_0260
% Mức: TH
\begin{ex}
% ID: AIQ_L10C6_FULL_0260
% Mức: TH
Ô tô qua khúc cua phẳng bán kính $42$ m với tốc độ $17$ m/s. Hệ số ma sát nghỉ tối thiểu, lấy $g=10$ m/s², bằng bao nhiêu?
\shortans{0,69}
\loigiai{
$\mu mg=mv^2/r\Rightarrow\mu=v^2/(gr)=17^2/(10\cdot42)$. Suy ra kết quả $0,69\ \mathrm{}$. Đáp án trả lời ngắn không ghi đơn vị.
}
\end{ex}

% ===== Câu 99 =====
% ID: AIQ_L10C6_FULL_0261
% Mức: VD
\begin{ex}
% ID: AIQ_L10C6_FULL_0261
% Mức: VD
Ô tô qua khúc cua phẳng bán kính $44$ m với tốc độ $18$ m/s. Hệ số ma sát nghỉ tối thiểu, lấy $g=10$ m/s², bằng bao nhiêu?
\shortans{0,74}
\loigiai{
$\mu mg=mv^2/r\Rightarrow\mu=v^2/(gr)=18^2/(10\cdot44)$. Suy ra kết quả $0,74\ \mathrm{}$. Đáp án trả lời ngắn không ghi đơn vị.
}
\end{ex}

% ===== Câu 100 =====
% ID: AIQ_L10C6_FULL_0262
% Mức: VD
\begin{ex}
% ID: AIQ_L10C6_FULL_0262
% Mức: VD
Ô tô qua khúc cua phẳng bán kính $46$ m với tốc độ $19$ m/s. Hệ số ma sát nghỉ tối thiểu, lấy $g=10$ m/s², bằng bao nhiêu?
\shortans{0,78}
\loigiai{
$\mu mg=mv^2/r\Rightarrow\mu=v^2/(gr)=19^2/(10\cdot46)$. Suy ra kết quả $0,78\ \mathrm{}$. Đáp án trả lời ngắn không ghi đơn vị.
}
\end{ex}

% ===== Câu 101 =====
% ID: AIQ_L10C6_FULL_0263
% Mức: VDC
\begin{ex}
% ID: AIQ_L10C6_FULL_0263
% Mức: VDC
Ô tô qua khúc cua phẳng bán kính $48$ m với tốc độ $20$ m/s. Hệ số ma sát nghỉ tối thiểu, lấy $g=10$ m/s², bằng bao nhiêu?
\shortans{0,83}
\loigiai{
$\mu mg=mv^2/r\Rightarrow\mu=v^2/(gr)=20^2/(10\cdot48)$. Suy ra kết quả $0,83\ \mathrm{}$. Đáp án trả lời ngắn không ghi đơn vị.
}
\end{ex}

% ===== Câu 102 =====
% ID: AIQ_L10C6_FULL_0264
% Mức: VDC
\begin{ex}
% ID: AIQ_L10C6_FULL_0264
% Mức: VDC
Ô tô qua khúc cua phẳng bán kính $50$ m với tốc độ $21$ m/s. Hệ số ma sát nghỉ tối thiểu, lấy $g=10$ m/s², bằng bao nhiêu?
\shortans{0,88}
\loigiai{
$\mu mg=mv^2/r\Rightarrow\mu=v^2/(gr)=21^2/(10\cdot50)$. Suy ra kết quả $0,88\ \mathrm{}$. Đáp án trả lời ngắn không ghi đơn vị.
}
\end{ex}

% ===== Câu 103 =====
% ID: AIQ_L10C6_FULL_0265
% Mức: TH
\begin{bt}
% ID: AIQ_L10C6_FULL_0265
% Mức: TH
Trình bày cách giải: Ô tô qua khúc cua phẳng bán kính $52$ m với tốc độ $22$ m/s. Hệ số ma sát nghỉ tối thiểu, lấy $g=10$ m/s², bằng bao nhiêu?
\loigiai{
$\mu mg=mv^2/r\Rightarrow\mu=v^2/(gr)=22^2/(10\cdot52)$. Suy ra kết quả $0,93\ \mathrm{}$.
}
\end{bt}

% ===== Câu 104 =====
% ID: AIQ_L10C6_FULL_0266
% Mức: TH
\begin{bt}
% ID: AIQ_L10C6_FULL_0266
% Mức: TH
Trình bày cách giải: Ô tô qua khúc cua phẳng bán kính $54$ m với tốc độ $23$ m/s. Hệ số ma sát nghỉ tối thiểu, lấy $g=10$ m/s², bằng bao nhiêu?
\loigiai{
$\mu mg=mv^2/r\Rightarrow\mu=v^2/(gr)=23^2/(10\cdot54)$. Suy ra kết quả $0,98\ \mathrm{}$.
}
\end{bt}

% ===== Câu 105 =====
% ID: AIQ_L10C6_FULL_0267
% Mức: VD
\begin{bt}
% ID: AIQ_L10C6_FULL_0267
% Mức: VD
Trình bày cách giải: Ô tô qua khúc cua phẳng bán kính $56$ m với tốc độ $24$ m/s. Hệ số ma sát nghỉ tối thiểu, lấy $g=10$ m/s², bằng bao nhiêu?
\loigiai{
$\mu mg=mv^2/r\Rightarrow\mu=v^2/(gr)=24^2/(10\cdot56)$. Suy ra kết quả $1,03\ \mathrm{}$.
}
\end{bt}

% ===== Câu 106 =====
% ID: AIQ_L10C6_FULL_0268
% Mức: VD
\begin{bt}
% ID: AIQ_L10C6_FULL_0268
% Mức: VD
Trình bày cách giải: Ô tô qua khúc cua phẳng bán kính $58$ m với tốc độ $25$ m/s. Hệ số ma sát nghỉ tối thiểu, lấy $g=10$ m/s², bằng bao nhiêu?
\loigiai{
$\mu mg=mv^2/r\Rightarrow\mu=v^2/(gr)=25^2/(10\cdot58)$. Suy ra kết quả $1,08\ \mathrm{}$.
}
\end{bt}

% ===== Câu 107 =====
% ID: AIQ_L10C6_FULL_0269
% Mức: VDC
\begin{bt}
% ID: AIQ_L10C6_FULL_0269
% Mức: VDC
Trình bày cách giải: Ô tô qua khúc cua phẳng bán kính $60$ m với tốc độ $26$ m/s. Hệ số ma sát nghỉ tối thiểu, lấy $g=10$ m/s², bằng bao nhiêu?
\loigiai{
$\mu mg=mv^2/r\Rightarrow\mu=v^2/(gr)=26^2/(10\cdot60)$. Suy ra kết quả $1,13\ \mathrm{}$.
}
\end{bt}

% ===== Câu 108 =====
% ID: AIQ_L10C6_FULL_0270
% Mức: VDC
\begin{bt}
% ID: AIQ_L10C6_FULL_0270
% Mức: VDC
Trình bày cách giải: Ô tô qua khúc cua phẳng bán kính $62$ m với tốc độ $27$ m/s. Hệ số ma sát nghỉ tối thiểu, lấy $g=10$ m/s², bằng bao nhiêu?
\loigiai{
$\mu mg=mv^2/r\Rightarrow\mu=v^2/(gr)=27^2/(10\cdot62)$. Suy ra kết quả $1,18\ \mathrm{}$.
}
\end{bt}

% ===== Câu 109 =====
% ID: AIQ_L10C6_FULL_0271
% Mức: NB
\dangbt{Chuyển động tròn trong mặt phẳng thẳng đứng}
\begin{ex}
% ID: AIQ_L10C6_FULL_0271
% Mức: NB
Vật khối lượng $0,2$ kg chuyển động qua đỉnh vòng tròn thẳng đứng bán kính $1$ m với tốc độ $4$ m/s. Tính lực ép của vòng lên vật, lấy $g=10$ m/s².
\choice
{\True $1,2\ \mathrm{N}$}
{$3,2\ \mathrm{N}$}
{$2\ \mathrm{N}$}
{$5,2\ \mathrm{N}$}
\loigiai{
Tại đỉnh: $N+mg=mv^2/r\Rightarrow N=mv^2/r-mg$. Suy ra kết quả $1,2\ \mathrm{N}$.
}
\end{ex}

% ===== Câu 110 =====
% ID: AIQ_L10C6_FULL_0272
% Mức: NB
\begin{ex}
% ID: AIQ_L10C6_FULL_0272
% Mức: NB
Vật khối lượng $0,3$ kg chuyển động qua đỉnh vòng tròn thẳng đứng bán kính $2$ m với tốc độ $5$ m/s. Tính lực ép của vòng lên vật, lấy $g=10$ m/s².
\choice
{$3,75\ \mathrm{N}$}
{\True $0,75\ \mathrm{N}$}
{$3\ \mathrm{N}$}
{$6,75\ \mathrm{N}$}
\loigiai{
Tại đỉnh: $N+mg=mv^2/r\Rightarrow N=mv^2/r-mg$. Suy ra kết quả $0,75\ \mathrm{N}$.
}
\end{ex}

% ===== Câu 111 =====
% ID: AIQ_L10C6_FULL_0273
% Mức: TH
\begin{ex}
% ID: AIQ_L10C6_FULL_0273
% Mức: TH
Vật khối lượng $0,4$ kg chuyển động qua đỉnh vòng tròn thẳng đứng bán kính $3$ m với tốc độ $6$ m/s. Tính lực ép của vòng lên vật, lấy $g=10$ m/s².
\choice
{$4,8\ \mathrm{N}$}
{$4\ \mathrm{N}$}
{\True $0,8\ \mathrm{N}$}
{$8,8\ \mathrm{N}$}
\loigiai{
Tại đỉnh: $N+mg=mv^2/r\Rightarrow N=mv^2/r-mg$. Suy ra kết quả $0,8\ \mathrm{N}$.
}
\end{ex}

% ===== Câu 112 =====
% ID: AIQ_L10C6_FULL_0274
% Mức: TH
\begin{ex}
% ID: AIQ_L10C6_FULL_0274
% Mức: TH
Vật khối lượng $0,5$ kg chuyển động qua đỉnh vòng tròn thẳng đứng bán kính $1$ m với tốc độ $7$ m/s. Tính lực ép của vòng lên vật, lấy $g=10$ m/s².
\choice
{$24,5\ \mathrm{N}$}
{$5\ \mathrm{N}$}
{$29,5\ \mathrm{N}$}
{\True $19,5\ \mathrm{N}$}
\loigiai{
Tại đỉnh: $N+mg=mv^2/r\Rightarrow N=mv^2/r-mg$. Suy ra kết quả $19,5\ \mathrm{N}$.
}
\end{ex}

% ===== Câu 113 =====
% ID: AIQ_L10C6_FULL_0275
% Mức: VD
\begin{ex}
% ID: AIQ_L10C6_FULL_0275
% Mức: VD
Vật khối lượng $0,2$ kg chuyển động qua đỉnh vòng tròn thẳng đứng bán kính $2$ m với tốc độ $8$ m/s. Tính lực ép của vòng lên vật, lấy $g=10$ m/s².
\choice
{\True $4,4\ \mathrm{N}$}
{$6,4\ \mathrm{N}$}
{$2\ \mathrm{N}$}
{$8,4\ \mathrm{N}$}
\loigiai{
Tại đỉnh: $N+mg=mv^2/r\Rightarrow N=mv^2/r-mg$. Suy ra kết quả $4,4\ \mathrm{N}$.
}
\end{ex}

% ===== Câu 114 =====
% ID: AIQ_L10C6_FULL_0276
% Mức: VD
\begin{ex}
% ID: AIQ_L10C6_FULL_0276
% Mức: VD
Vật khối lượng $0,3$ kg chuyển động qua đỉnh vòng tròn thẳng đứng bán kính $3$ m với tốc độ $9$ m/s. Tính lực ép của vòng lên vật, lấy $g=10$ m/s².
\choice
{$8,1\ \mathrm{N}$}
{\True $5,1\ \mathrm{N}$}
{$3\ \mathrm{N}$}
{$11,1\ \mathrm{N}$}
\loigiai{
Tại đỉnh: $N+mg=mv^2/r\Rightarrow N=mv^2/r-mg$. Suy ra kết quả $5,1\ \mathrm{N}$.
}
\end{ex}

% ===== Câu 115 =====
% ID: AIQ_L10C6_FULL_0277
% Mức: VD
\begin{ex}
% ID: AIQ_L10C6_FULL_0277
% Mức: VD
Vật khối lượng $0,4$ kg chuyển động qua đỉnh vòng tròn thẳng đứng bán kính $1$ m với tốc độ $10$ m/s. Tính lực ép của vòng lên vật, lấy $g=10$ m/s².
\choice
{$40\ \mathrm{N}$}
{$4\ \mathrm{N}$}
{\True $36\ \mathrm{N}$}
{$44\ \mathrm{N}$}
\loigiai{
Tại đỉnh: $N+mg=mv^2/r\Rightarrow N=mv^2/r-mg$. Suy ra kết quả $36\ \mathrm{N}$.
}
\end{ex}

% ===== Câu 116 =====
% ID: AIQ_L10C6_FULL_0278
% Mức: VDC
\begin{ex}
% ID: AIQ_L10C6_FULL_0278
% Mức: VDC
Vật khối lượng $0,5$ kg chuyển động qua đỉnh vòng tròn thẳng đứng bán kính $2$ m với tốc độ $11$ m/s. Tính lực ép của vòng lên vật, lấy $g=10$ m/s².
\choice
{$30,25\ \mathrm{N}$}
{$5\ \mathrm{N}$}
{$35,25\ \mathrm{N}$}
{\True $25,25\ \mathrm{N}$}
\loigiai{
Tại đỉnh: $N+mg=mv^2/r\Rightarrow N=mv^2/r-mg$. Suy ra kết quả $25,25\ \mathrm{N}$.
}
\end{ex}

% ===== Câu 117 =====
% ID: AIQ_L10C6_FULL_0279
% Mức: VDC
\begin{ex}
% ID: AIQ_L10C6_FULL_0279
% Mức: VDC
Vật khối lượng $0,2$ kg chuyển động qua đỉnh vòng tròn thẳng đứng bán kính $3$ m với tốc độ $12$ m/s. Tính lực ép của vòng lên vật, lấy $g=10$ m/s².
\choice
{\True $7,6\ \mathrm{N}$}
{$9,6\ \mathrm{N}$}
{$2\ \mathrm{N}$}
{$11,6\ \mathrm{N}$}
\loigiai{
Tại đỉnh: $N+mg=mv^2/r\Rightarrow N=mv^2/r-mg$. Suy ra kết quả $7,6\ \mathrm{N}$.
}
\end{ex}

% ===== Câu 118 =====
% ID: AIQ_L10C6_FULL_0280
% Mức: VDC
\begin{ex}
% ID: AIQ_L10C6_FULL_0280
% Mức: VDC
Vật khối lượng $0,3$ kg chuyển động qua đỉnh vòng tròn thẳng đứng bán kính $1$ m với tốc độ $13$ m/s. Tính lực ép của vòng lên vật, lấy $g=10$ m/s².
\choice
{$50,7\ \mathrm{N}$}
{\True $47,7\ \mathrm{N}$}
{$3\ \mathrm{N}$}
{$53,7\ \mathrm{N}$}
\loigiai{
Tại đỉnh: $N+mg=mv^2/r\Rightarrow N=mv^2/r-mg$. Suy ra kết quả $47,7\ \mathrm{N}$.
}
\end{ex}

% ===== Câu 119 =====
% ID: AIQ_L10C6_FULL_0281
% Mức: NB
\begin{ex}
% ID: AIQ_L10C6_FULL_0281
% Mức: NB
Xét các phát biểu thuộc dạng “Chuyển động tròn trong mặt phẳng thẳng đứng” ở tình huống số 1.
\choiceTF
{\True Tại đỉnh vòng tròn, gia tốc hướng về tâm.}
{\True Trọng lực có thể góp phần tạo lực hướng tâm.}
{Lực ép tại đỉnh luôn bằng $mg$.}
{\True Điều kiện không rời quỹ đạo tại đỉnh là $v^2/r\ge g$ trong trường hợp vật ở phía trong vòng.}
\loigiai{
\begin{itemchoice}
\item (Đúng) Tại đỉnh vòng tròn, gia tốc hướng về tâm. (Vì): Phù hợp với định nghĩa hoặc hệ thức vật lí. b) Đúng: Phù hợp với định nghĩa hoặc hệ thức vật lí. c) Sai: Không phù hợp với định nghĩa hoặc hệ thức vật lí. d) Đúng: Phù hợp với định nghĩa hoặc hệ thức vật lí
\item (Đúng) Trọng lực có thể góp phần tạo lực hướng tâm.
\item (Sai) Lực ép tại đỉnh luôn bằng $mg$.
\item (Đúng) Điều kiện không rời quỹ đạo tại đỉnh là $v^2/r\ge g$ trong trường hợp vật ở phía trong vòng.
\end{itemchoice}
}
\end{ex}

% ===== Câu 120 =====
% ID: AIQ_L10C6_FULL_0282
% Mức: TH
\begin{ex}
% ID: AIQ_L10C6_FULL_0282
% Mức: TH
Xét các phát biểu thuộc dạng “Chuyển động tròn trong mặt phẳng thẳng đứng” ở tình huống số 2.
\choiceTF
{\True Tại đỉnh vòng tròn, gia tốc hướng về tâm.}
{\True Trọng lực có thể góp phần tạo lực hướng tâm.}
{Lực ép tại đỉnh luôn bằng $mg$.}
{\True Điều kiện không rời quỹ đạo tại đỉnh là $v^2/r\ge g$ trong trường hợp vật ở phía trong vòng.}
\loigiai{
\begin{itemchoice}
\item (Đúng) Tại đỉnh vòng tròn, gia tốc hướng về tâm. (Vì): Phù hợp với định nghĩa hoặc hệ thức vật lí. b) Đúng: Phù hợp với định nghĩa hoặc hệ thức vật lí. c) Sai: Không phù hợp với định nghĩa hoặc hệ thức vật lí. d) Đúng: Phù hợp với định nghĩa hoặc hệ thức vật lí
\item (Đúng) Trọng lực có thể góp phần tạo lực hướng tâm.
\item (Sai) Lực ép tại đỉnh luôn bằng $mg$.
\item (Đúng) Điều kiện không rời quỹ đạo tại đỉnh là $v^2/r\ge g$ trong trường hợp vật ở phía trong vòng.
\end{itemchoice}
}
\end{ex}

% ===== Câu 121 =====
% ID: AIQ_L10C6_FULL_0283
% Mức: TH
\begin{ex}
% ID: AIQ_L10C6_FULL_0283
% Mức: TH
Xét các phát biểu thuộc dạng “Chuyển động tròn trong mặt phẳng thẳng đứng” ở tình huống số 3.
\choiceTF
{\True Tại đỉnh vòng tròn, gia tốc hướng về tâm.}
{\True Trọng lực có thể góp phần tạo lực hướng tâm.}
{Lực ép tại đỉnh luôn bằng $mg$.}
{\True Điều kiện không rời quỹ đạo tại đỉnh là $v^2/r\ge g$ trong trường hợp vật ở phía trong vòng.}
\loigiai{
\begin{itemchoice}
\item (Đúng) Tại đỉnh vòng tròn, gia tốc hướng về tâm. (Vì): Phù hợp với định nghĩa hoặc hệ thức vật lí. b) Đúng: Phù hợp với định nghĩa hoặc hệ thức vật lí. c) Sai: Không phù hợp với định nghĩa hoặc hệ thức vật lí. d) Đúng: Phù hợp với định nghĩa hoặc hệ thức vật lí
\item (Đúng) Trọng lực có thể góp phần tạo lực hướng tâm.
\item (Sai) Lực ép tại đỉnh luôn bằng $mg$.
\item (Đúng) Điều kiện không rời quỹ đạo tại đỉnh là $v^2/r\ge g$ trong trường hợp vật ở phía trong vòng.
\end{itemchoice}
}
\end{ex}

% ===== Câu 122 =====
% ID: AIQ_L10C6_FULL_0284
% Mức: VD
\begin{ex}
% ID: AIQ_L10C6_FULL_0284
% Mức: VD
Xét các phát biểu thuộc dạng “Chuyển động tròn trong mặt phẳng thẳng đứng” ở tình huống số 4.
\choiceTF
{\True Tại đỉnh vòng tròn, gia tốc hướng về tâm.}
{\True Trọng lực có thể góp phần tạo lực hướng tâm.}
{Lực ép tại đỉnh luôn bằng $mg$.}
{\True Điều kiện không rời quỹ đạo tại đỉnh là $v^2/r\ge g$ trong trường hợp vật ở phía trong vòng.}
\loigiai{
\begin{itemchoice}
\item (Đúng) Tại đỉnh vòng tròn, gia tốc hướng về tâm. (Vì): Phù hợp với định nghĩa hoặc hệ thức vật lí. b) Đúng: Phù hợp với định nghĩa hoặc hệ thức vật lí. c) Sai: Không phù hợp với định nghĩa hoặc hệ thức vật lí. d) Đúng: Phù hợp với định nghĩa hoặc hệ thức vật lí
\item (Đúng) Trọng lực có thể góp phần tạo lực hướng tâm.
\item (Sai) Lực ép tại đỉnh luôn bằng $mg$.
\item (Đúng) Điều kiện không rời quỹ đạo tại đỉnh là $v^2/r\ge g$ trong trường hợp vật ở phía trong vòng.
\end{itemchoice}
}
\end{ex}

% ===== Câu 123 =====
% ID: AIQ_L10C6_FULL_0285
% Mức: VDC
\begin{ex}
% ID: AIQ_L10C6_FULL_0285
% Mức: VDC
Xét các phát biểu thuộc dạng “Chuyển động tròn trong mặt phẳng thẳng đứng” ở tình huống số 5.
\choiceTF
{\True Tại đỉnh vòng tròn, gia tốc hướng về tâm.}
{\True Trọng lực có thể góp phần tạo lực hướng tâm.}
{Lực ép tại đỉnh luôn bằng $mg$.}
{\True Điều kiện không rời quỹ đạo tại đỉnh là $v^2/r\ge g$ trong trường hợp vật ở phía trong vòng.}
\loigiai{
\begin{itemchoice}
\item (Đúng) Tại đỉnh vòng tròn, gia tốc hướng về tâm. (Vì): Phù hợp với định nghĩa hoặc hệ thức vật lí. b) Đúng: Phù hợp với định nghĩa hoặc hệ thức vật lí. c) Sai: Không phù hợp với định nghĩa hoặc hệ thức vật lí. d) Đúng: Phù hợp với định nghĩa hoặc hệ thức vật lí
\item (Đúng) Trọng lực có thể góp phần tạo lực hướng tâm.
\item (Sai) Lực ép tại đỉnh luôn bằng $mg$.
\item (Đúng) Điều kiện không rời quỹ đạo tại đỉnh là $v^2/r\ge g$ trong trường hợp vật ở phía trong vòng.
\end{itemchoice}
}
\end{ex}

% ===== Câu 124 =====
% ID: AIQ_L10C6_FULL_0286
% Mức: TH
\begin{ex}
% ID: AIQ_L10C6_FULL_0286
% Mức: TH
Vật khối lượng $0,4$ kg chuyển động qua đỉnh vòng tròn thẳng đứng bán kính $2$ m với tốc độ $14$ m/s. Tính lực ép của vòng lên vật, lấy $g=10$ m/s².
\shortans{35,2}
\loigiai{
Tại đỉnh: $N+mg=mv^2/r\Rightarrow N=mv^2/r-mg$. Suy ra kết quả $35,2\ \mathrm{N}$. Đáp án trả lời ngắn không ghi đơn vị.
}
\end{ex}

% ===== Câu 125 =====
% ID: AIQ_L10C6_FULL_0287
% Mức: TH
\begin{ex}
% ID: AIQ_L10C6_FULL_0287
% Mức: TH
Vật khối lượng $0,5$ kg chuyển động qua đỉnh vòng tròn thẳng đứng bán kính $3$ m với tốc độ $15$ m/s. Tính lực ép của vòng lên vật, lấy $g=10$ m/s².
\shortans{32,5}
\loigiai{
Tại đỉnh: $N+mg=mv^2/r\Rightarrow N=mv^2/r-mg$. Suy ra kết quả $32,5\ \mathrm{N}$. Đáp án trả lời ngắn không ghi đơn vị.
}
\end{ex}

% ===== Câu 126 =====
% ID: AIQ_L10C6_FULL_0288
% Mức: VD
\begin{ex}
% ID: AIQ_L10C6_FULL_0288
% Mức: VD
Vật khối lượng $0,2$ kg chuyển động qua đỉnh vòng tròn thẳng đứng bán kính $1$ m với tốc độ $16$ m/s. Tính lực ép của vòng lên vật, lấy $g=10$ m/s².
\shortans{49,2}
\loigiai{
Tại đỉnh: $N+mg=mv^2/r\Rightarrow N=mv^2/r-mg$. Suy ra kết quả $49,2\ \mathrm{N}$. Đáp án trả lời ngắn không ghi đơn vị.
}
\end{ex}

% ===== Câu 127 =====
% ID: AIQ_L10C6_FULL_0289
% Mức: VD
\begin{ex}
% ID: AIQ_L10C6_FULL_0289
% Mức: VD
Vật khối lượng $0,3$ kg chuyển động qua đỉnh vòng tròn thẳng đứng bán kính $2$ m với tốc độ $17$ m/s. Tính lực ép của vòng lên vật, lấy $g=10$ m/s².
\shortans{40,35}
\loigiai{
Tại đỉnh: $N+mg=mv^2/r\Rightarrow N=mv^2/r-mg$. Suy ra kết quả $40,35\ \mathrm{N}$. Đáp án trả lời ngắn không ghi đơn vị.
}
\end{ex}

% ===== Câu 128 =====
% ID: AIQ_L10C6_FULL_0290
% Mức: VDC
\begin{ex}
% ID: AIQ_L10C6_FULL_0290
% Mức: VDC
Vật khối lượng $0,4$ kg chuyển động qua đỉnh vòng tròn thẳng đứng bán kính $3$ m với tốc độ $18$ m/s. Tính lực ép của vòng lên vật, lấy $g=10$ m/s².
\shortans{39,2}
\loigiai{
Tại đỉnh: $N+mg=mv^2/r\Rightarrow N=mv^2/r-mg$. Suy ra kết quả $39,2\ \mathrm{N}$. Đáp án trả lời ngắn không ghi đơn vị.
}
\end{ex}

% ===== Câu 129 =====
% ID: AIQ_L10C6_FULL_0291
% Mức: VDC
\begin{ex}
% ID: AIQ_L10C6_FULL_0291
% Mức: VDC
Vật khối lượng $0,5$ kg chuyển động qua đỉnh vòng tròn thẳng đứng bán kính $1$ m với tốc độ $19$ m/s. Tính lực ép của vòng lên vật, lấy $g=10$ m/s².
\shortans{175,5}
\loigiai{
Tại đỉnh: $N+mg=mv^2/r\Rightarrow N=mv^2/r-mg$. Suy ra kết quả $175,5\ \mathrm{N}$. Đáp án trả lời ngắn không ghi đơn vị.
}
\end{ex}

% ===== Câu 130 =====
% ID: AIQ_L10C6_FULL_0292
% Mức: TH
\begin{bt}
% ID: AIQ_L10C6_FULL_0292
% Mức: TH
Trình bày cách giải: Vật khối lượng $0,2$ kg chuyển động qua đỉnh vòng tròn thẳng đứng bán kính $2$ m với tốc độ $20$ m/s. Tính lực ép của vòng lên vật, lấy $g=10$ m/s².
\loigiai{
Tại đỉnh: $N+mg=mv^2/r\Rightarrow N=mv^2/r-mg$. Suy ra kết quả $38\ \mathrm{N}$.
}
\end{bt}

% ===== Câu 131 =====
% ID: AIQ_L10C6_FULL_0293
% Mức: TH
\begin{bt}
% ID: AIQ_L10C6_FULL_0293
% Mức: TH
Trình bày cách giải: Vật khối lượng $0,3$ kg chuyển động qua đỉnh vòng tròn thẳng đứng bán kính $3$ m với tốc độ $21$ m/s. Tính lực ép của vòng lên vật, lấy $g=10$ m/s².
\loigiai{
Tại đỉnh: $N+mg=mv^2/r\Rightarrow N=mv^2/r-mg$. Suy ra kết quả $41,1\ \mathrm{N}$.
}
\end{bt}

% ===== Câu 132 =====
% ID: AIQ_L10C6_FULL_0294
% Mức: VD
\begin{bt}
% ID: AIQ_L10C6_FULL_0294
% Mức: VD
Trình bày cách giải: Vật khối lượng $0,4$ kg chuyển động qua đỉnh vòng tròn thẳng đứng bán kính $1$ m với tốc độ $22$ m/s. Tính lực ép của vòng lên vật, lấy $g=10$ m/s².
\loigiai{
Tại đỉnh: $N+mg=mv^2/r\Rightarrow N=mv^2/r-mg$. Suy ra kết quả $189,6\ \mathrm{N}$.
}
\end{bt}

% ===== Câu 133 =====
% ID: AIQ_L10C6_FULL_0295
% Mức: VD
\begin{bt}
% ID: AIQ_L10C6_FULL_0295
% Mức: VD
Trình bày cách giải: Vật khối lượng $0,5$ kg chuyển động qua đỉnh vòng tròn thẳng đứng bán kính $2$ m với tốc độ $23$ m/s. Tính lực ép của vòng lên vật, lấy $g=10$ m/s².
\loigiai{
Tại đỉnh: $N+mg=mv^2/r\Rightarrow N=mv^2/r-mg$. Suy ra kết quả $127,25\ \mathrm{N}$.
}
\end{bt}

% ===== Câu 134 =====
% ID: AIQ_L10C6_FULL_0296
% Mức: VDC
\begin{bt}
% ID: AIQ_L10C6_FULL_0296
% Mức: VDC
Trình bày cách giải: Vật khối lượng $0,2$ kg chuyển động qua đỉnh vòng tròn thẳng đứng bán kính $3$ m với tốc độ $24$ m/s. Tính lực ép của vòng lên vật, lấy $g=10$ m/s².
\loigiai{
Tại đỉnh: $N+mg=mv^2/r\Rightarrow N=mv^2/r-mg$. Suy ra kết quả $36,4\ \mathrm{N}$.
}
\end{bt}

% ===== Câu 135 =====
% ID: AIQ_L10C6_FULL_0297
% Mức: VDC
\begin{bt}
% ID: AIQ_L10C6_FULL_0297
% Mức: VDC
Trình bày cách giải: Vật khối lượng $0,3$ kg chuyển động qua đỉnh vòng tròn thẳng đứng bán kính $1$ m với tốc độ $25$ m/s. Tính lực ép của vòng lên vật, lấy $g=10$ m/s².
\loigiai{
Tại đỉnh: $N+mg=mv^2/r\Rightarrow N=mv^2/r-mg$. Suy ra kết quả $184,5\ \mathrm{N}$.
}
\end{bt}

% ===== Câu 136 =====
% ID: AIQ_L10C6_FULL_0298
% Mức: NB
\dangbt{Bài toán tổng hợp lực hướng tâm gắn với thực tiễn}
\begin{ex}
% ID: AIQ_L10C6_FULL_0298
% Mức: NB
Một vật khối lượng $0,5$ kg quay tròn đều bán kính $2$ m với tốc độ $5$ m/s. Tính lực hướng tâm cần thiết.
\choice
{\True $6,25\ \mathrm{N}$}
{$1,25\ \mathrm{N}$}
{$25\ \mathrm{N}$}
{$12,5\ \mathrm{N}$}
\loigiai{
$F_{ht}=mv^2/r=0,5\cdot5^2/2$. Suy ra kết quả $6,25\ \mathrm{N}$.
}
\end{ex}

% ===== Câu 137 =====
% ID: AIQ_L10C6_FULL_0299
% Mức: NB
\begin{ex}
% ID: AIQ_L10C6_FULL_0299
% Mức: NB
Một vật khối lượng $0,6$ kg quay tròn đều bán kính $3$ m với tốc độ $6$ m/s. Tính lực hướng tâm cần thiết.
\choice
{$1,2\ \mathrm{N}$}
{\True $7,2\ \mathrm{N}$}
{$64,8\ \mathrm{N}$}
{$12\ \mathrm{N}$}
\loigiai{
$F_{ht}=mv^2/r=0,6\cdot6^2/3$. Suy ra kết quả $7,2\ \mathrm{N}$.
}
\end{ex}

% ===== Câu 138 =====
% ID: AIQ_L10C6_FULL_0300
% Mức: TH
\begin{ex}
% ID: AIQ_L10C6_FULL_0300
% Mức: TH
Một vật khối lượng $0,7$ kg quay tròn đều bán kính $4$ m với tốc độ $7$ m/s. Tính lực hướng tâm cần thiết.
\choice
{$1,22\ \mathrm{N}$}
{$137,2\ \mathrm{N}$}
{\True $8,57\ \mathrm{N}$}
{$12,25\ \mathrm{N}$}
\loigiai{
$F_{ht}=mv^2/r=0,7\cdot7^2/4$. Suy ra kết quả $8,57\ \mathrm{N}$.
}
\end{ex}

% ===== Câu 139 =====
% ID: AIQ_L10C6_FULL_0301
% Mức: TH
\begin{ex}
% ID: AIQ_L10C6_FULL_0301
% Mức: TH
Một vật khối lượng $0,8$ kg quay tròn đều bán kính $5$ m với tốc độ $8$ m/s. Tính lực hướng tâm cần thiết.
\choice
{$1,28\ \mathrm{N}$}
{$256\ \mathrm{N}$}
{$12,8\ \mathrm{N}$}
{\True $10,24\ \mathrm{N}$}
\loigiai{
$F_{ht}=mv^2/r=0,8\cdot8^2/5$. Suy ra kết quả $10,24\ \mathrm{N}$.
}
\end{ex}

% ===== Câu 140 =====
% ID: AIQ_L10C6_FULL_0302
% Mức: VD
\begin{ex}
% ID: AIQ_L10C6_FULL_0302
% Mức: VD
Một vật khối lượng $0,9$ kg quay tròn đều bán kính $2$ m với tốc độ $9$ m/s. Tính lực hướng tâm cần thiết.
\choice
{\True $36,45\ \mathrm{N}$}
{$4,05\ \mathrm{N}$}
{$145,8\ \mathrm{N}$}
{$40,5\ \mathrm{N}$}
\loigiai{
$F_{ht}=mv^2/r=0,9\cdot9^2/2$. Suy ra kết quả $36,45\ \mathrm{N}$.
}
\end{ex}

% ===== Câu 141 =====
% ID: AIQ_L10C6_FULL_0303
% Mức: VD
\begin{ex}
% ID: AIQ_L10C6_FULL_0303
% Mức: VD
Một vật khối lượng $0,5$ kg quay tròn đều bán kính $3$ m với tốc độ $10$ m/s. Tính lực hướng tâm cần thiết.
\choice
{$1,67\ \mathrm{N}$}
{\True $16,67\ \mathrm{N}$}
{$150\ \mathrm{N}$}
{$33,33\ \mathrm{N}$}
\loigiai{
$F_{ht}=mv^2/r=0,5\cdot10^2/3$. Suy ra kết quả $16,67\ \mathrm{N}$.
}
\end{ex}

% ===== Câu 142 =====
% ID: AIQ_L10C6_FULL_0304
% Mức: VD
\begin{ex}
% ID: AIQ_L10C6_FULL_0304
% Mức: VD
Một vật khối lượng $0,6$ kg quay tròn đều bán kính $4$ m với tốc độ $11$ m/s. Tính lực hướng tâm cần thiết.
\choice
{$1,65\ \mathrm{N}$}
{$290,4\ \mathrm{N}$}
{\True $18,15\ \mathrm{N}$}
{$30,25\ \mathrm{N}$}
\loigiai{
$F_{ht}=mv^2/r=0,6\cdot11^2/4$. Suy ra kết quả $18,15\ \mathrm{N}$.
}
\end{ex}

% ===== Câu 143 =====
% ID: AIQ_L10C6_FULL_0305
% Mức: VDC
\begin{ex}
% ID: AIQ_L10C6_FULL_0305
% Mức: VDC
Một vật khối lượng $0,7$ kg quay tròn đều bán kính $5$ m với tốc độ $12$ m/s. Tính lực hướng tâm cần thiết.
\choice
{$1,68\ \mathrm{N}$}
{$504\ \mathrm{N}$}
{$28,8\ \mathrm{N}$}
{\True $20,16\ \mathrm{N}$}
\loigiai{
$F_{ht}=mv^2/r=0,7\cdot12^2/5$. Suy ra kết quả $20,16\ \mathrm{N}$.
}
\end{ex}

% ===== Câu 144 =====
% ID: AIQ_L10C6_FULL_0306
% Mức: VDC
\begin{ex}
% ID: AIQ_L10C6_FULL_0306
% Mức: VDC
Một vật khối lượng $0,8$ kg quay tròn đều bán kính $2$ m với tốc độ $13$ m/s. Tính lực hướng tâm cần thiết.
\choice
{\True $67,6\ \mathrm{N}$}
{$5,2\ \mathrm{N}$}
{$270,4\ \mathrm{N}$}
{$84,5\ \mathrm{N}$}
\loigiai{
$F_{ht}=mv^2/r=0,8\cdot13^2/2$. Suy ra kết quả $67,6\ \mathrm{N}$.
}
\end{ex}

% ===== Câu 145 =====
% ID: AIQ_L10C6_FULL_0307
% Mức: VDC
\begin{ex}
% ID: AIQ_L10C6_FULL_0307
% Mức: VDC
Một vật khối lượng $0,9$ kg quay tròn đều bán kính $3$ m với tốc độ $14$ m/s. Tính lực hướng tâm cần thiết.
\choice
{$4,2\ \mathrm{N}$}
{\True $58,8\ \mathrm{N}$}
{$529,2\ \mathrm{N}$}
{$65,33\ \mathrm{N}$}
\loigiai{
$F_{ht}=mv^2/r=0,9\cdot14^2/3$. Suy ra kết quả $58,8\ \mathrm{N}$.
}
\end{ex}

% ===== Câu 146 =====
% ID: AIQ_L10C6_FULL_0308
% Mức: NB
\begin{ex}
% ID: AIQ_L10C6_FULL_0308
% Mức: NB
Xét các phát biểu thuộc dạng “Bài toán tổng hợp lực hướng tâm gắn với thực tiễn” ở tình huống số 1.
\choiceTF
{\True Vệ tinh chuyển động tròn cần lực hướng tâm.}
{\True Trọng lực có thể là lực hướng tâm.}
{Lực hướng tâm sinh công trong chuyển động tròn đều.}
{\True Thiết kế đường cong phải xét bán kính và tốc độ.}
\loigiai{
\begin{itemchoice}
\item (Đúng) Vệ tinh chuyển động tròn cần lực hướng tâm. (Vì): Phù hợp với định nghĩa hoặc hệ thức vật lí. b) Đúng: Phù hợp với định nghĩa hoặc hệ thức vật lí. c) Sai: Không phù hợp với định nghĩa hoặc hệ thức vật lí. d) Đúng: Phù hợp với định nghĩa hoặc hệ thức vật lí
\item (Đúng) Trọng lực có thể là lực hướng tâm.
\item (Sai) Lực hướng tâm sinh công trong chuyển động tròn đều.
\item (Đúng) Thiết kế đường cong phải xét bán kính và tốc độ.
\end{itemchoice}
}
\end{ex}

% ===== Câu 147 =====
% ID: AIQ_L10C6_FULL_0309
% Mức: TH
\begin{ex}
% ID: AIQ_L10C6_FULL_0309
% Mức: TH
Xét các phát biểu thuộc dạng “Bài toán tổng hợp lực hướng tâm gắn với thực tiễn” ở tình huống số 2.
\choiceTF
{\True Vệ tinh chuyển động tròn cần lực hướng tâm.}
{\True Trọng lực có thể là lực hướng tâm.}
{Lực hướng tâm sinh công trong chuyển động tròn đều.}
{\True Thiết kế đường cong phải xét bán kính và tốc độ.}
\loigiai{
\begin{itemchoice}
\item (Đúng) Vệ tinh chuyển động tròn cần lực hướng tâm. (Vì): Phù hợp với định nghĩa hoặc hệ thức vật lí. b) Đúng: Phù hợp với định nghĩa hoặc hệ thức vật lí. c) Sai: Không phù hợp với định nghĩa hoặc hệ thức vật lí. d) Đúng: Phù hợp với định nghĩa hoặc hệ thức vật lí
\item (Đúng) Trọng lực có thể là lực hướng tâm.
\item (Sai) Lực hướng tâm sinh công trong chuyển động tròn đều.
\item (Đúng) Thiết kế đường cong phải xét bán kính và tốc độ.
\end{itemchoice}
}
\end{ex}

% ===== Câu 148 =====
% ID: AIQ_L10C6_FULL_0310
% Mức: TH
\begin{ex}
% ID: AIQ_L10C6_FULL_0310
% Mức: TH
Xét các phát biểu thuộc dạng “Bài toán tổng hợp lực hướng tâm gắn với thực tiễn” ở tình huống số 3.
\choiceTF
{\True Vệ tinh chuyển động tròn cần lực hướng tâm.}
{\True Trọng lực có thể là lực hướng tâm.}
{Lực hướng tâm sinh công trong chuyển động tròn đều.}
{\True Thiết kế đường cong phải xét bán kính và tốc độ.}
\loigiai{
\begin{itemchoice}
\item (Đúng) Vệ tinh chuyển động tròn cần lực hướng tâm. (Vì): Phù hợp với định nghĩa hoặc hệ thức vật lí. b) Đúng: Phù hợp với định nghĩa hoặc hệ thức vật lí. c) Sai: Không phù hợp với định nghĩa hoặc hệ thức vật lí. d) Đúng: Phù hợp với định nghĩa hoặc hệ thức vật lí
\item (Đúng) Trọng lực có thể là lực hướng tâm.
\item (Sai) Lực hướng tâm sinh công trong chuyển động tròn đều.
\item (Đúng) Thiết kế đường cong phải xét bán kính và tốc độ.
\end{itemchoice}
}
\end{ex}

% ===== Câu 149 =====
% ID: AIQ_L10C6_FULL_0311
% Mức: VD
\begin{ex}
% ID: AIQ_L10C6_FULL_0311
% Mức: VD
Xét các phát biểu thuộc dạng “Bài toán tổng hợp lực hướng tâm gắn với thực tiễn” ở tình huống số 4.
\choiceTF
{\True Vệ tinh chuyển động tròn cần lực hướng tâm.}
{\True Trọng lực có thể là lực hướng tâm.}
{Lực hướng tâm sinh công trong chuyển động tròn đều.}
{\True Thiết kế đường cong phải xét bán kính và tốc độ.}
\loigiai{
\begin{itemchoice}
\item (Đúng) Vệ tinh chuyển động tròn cần lực hướng tâm. (Vì): Phù hợp với định nghĩa hoặc hệ thức vật lí. b) Đúng: Phù hợp với định nghĩa hoặc hệ thức vật lí. c) Sai: Không phù hợp với định nghĩa hoặc hệ thức vật lí. d) Đúng: Phù hợp với định nghĩa hoặc hệ thức vật lí
\item (Đúng) Trọng lực có thể là lực hướng tâm.
\item (Sai) Lực hướng tâm sinh công trong chuyển động tròn đều.
\item (Đúng) Thiết kế đường cong phải xét bán kính và tốc độ.
\end{itemchoice}
}
\end{ex}

% ===== Câu 150 =====
% ID: AIQ_L10C6_FULL_0312
% Mức: VDC
\begin{ex}
% ID: AIQ_L10C6_FULL_0312
% Mức: VDC
Xét các phát biểu thuộc dạng “Bài toán tổng hợp lực hướng tâm gắn với thực tiễn” ở tình huống số 5.
\choiceTF
{\True Vệ tinh chuyển động tròn cần lực hướng tâm.}
{\True Trọng lực có thể là lực hướng tâm.}
{Lực hướng tâm sinh công trong chuyển động tròn đều.}
{\True Thiết kế đường cong phải xét bán kính và tốc độ.}
\loigiai{
\begin{itemchoice}
\item (Đúng) Vệ tinh chuyển động tròn cần lực hướng tâm. (Vì): Phù hợp với định nghĩa hoặc hệ thức vật lí. b) Đúng: Phù hợp với định nghĩa hoặc hệ thức vật lí. c) Sai: Không phù hợp với định nghĩa hoặc hệ thức vật lí. d) Đúng: Phù hợp với định nghĩa hoặc hệ thức vật lí
\item (Đúng) Trọng lực có thể là lực hướng tâm.
\item (Sai) Lực hướng tâm sinh công trong chuyển động tròn đều.
\item (Đúng) Thiết kế đường cong phải xét bán kính và tốc độ.
\end{itemchoice}
}
\end{ex}

% ===== Câu 151 =====
% ID: AIQ_L10C6_FULL_0313
% Mức: TH
\begin{ex}
% ID: AIQ_L10C6_FULL_0313
% Mức: TH
Một vật khối lượng $0,5$ kg quay tròn đều bán kính $4$ m với tốc độ $15$ m/s. Tính lực hướng tâm cần thiết.
\shortans{28,13}
\loigiai{
$F_{ht}=mv^2/r=0,5\cdot15^2/4$. Suy ra kết quả $28,13\ \mathrm{N}$. Đáp án trả lời ngắn không ghi đơn vị.
}
\end{ex}

% ===== Câu 152 =====
% ID: AIQ_L10C6_FULL_0314
% Mức: TH
\begin{ex}
% ID: AIQ_L10C6_FULL_0314
% Mức: TH
Một vật khối lượng $0,6$ kg quay tròn đều bán kính $5$ m với tốc độ $16$ m/s. Tính lực hướng tâm cần thiết.
\shortans{30,72}
\loigiai{
$F_{ht}=mv^2/r=0,6\cdot16^2/5$. Suy ra kết quả $30,72\ \mathrm{N}$. Đáp án trả lời ngắn không ghi đơn vị.
}
\end{ex}

% ===== Câu 153 =====
% ID: AIQ_L10C6_FULL_0315
% Mức: VD
\begin{ex}
% ID: AIQ_L10C6_FULL_0315
% Mức: VD
Một vật khối lượng $0,7$ kg quay tròn đều bán kính $2$ m với tốc độ $17$ m/s. Tính lực hướng tâm cần thiết.
\shortans{101,15}
\loigiai{
$F_{ht}=mv^2/r=0,7\cdot17^2/2$. Suy ra kết quả $101,15\ \mathrm{N}$. Đáp án trả lời ngắn không ghi đơn vị.
}
\end{ex}

% ===== Câu 154 =====
% ID: AIQ_L10C6_FULL_0316
% Mức: VD
\begin{ex}
% ID: AIQ_L10C6_FULL_0316
% Mức: VD
Một vật khối lượng $0,8$ kg quay tròn đều bán kính $3$ m với tốc độ $18$ m/s. Tính lực hướng tâm cần thiết.
\shortans{86,4}
\loigiai{
$F_{ht}=mv^2/r=0,8\cdot18^2/3$. Suy ra kết quả $86,4\ \mathrm{N}$. Đáp án trả lời ngắn không ghi đơn vị.
}
\end{ex}

% ===== Câu 155 =====
% ID: AIQ_L10C6_FULL_0317
% Mức: VDC
\begin{ex}
% ID: AIQ_L10C6_FULL_0317
% Mức: VDC
Một vật khối lượng $0,9$ kg quay tròn đều bán kính $4$ m với tốc độ $19$ m/s. Tính lực hướng tâm cần thiết.
\shortans{81,23}
\loigiai{
$F_{ht}=mv^2/r=0,9\cdot19^2/4$. Suy ra kết quả $81,23\ \mathrm{N}$. Đáp án trả lời ngắn không ghi đơn vị.
}
\end{ex}

% ===== Câu 156 =====
% ID: AIQ_L10C6_FULL_0318
% Mức: VDC
\begin{ex}
% ID: AIQ_L10C6_FULL_0318
% Mức: VDC
Một vật khối lượng $0,5$ kg quay tròn đều bán kính $5$ m với tốc độ $20$ m/s. Tính lực hướng tâm cần thiết.
\shortans{40}
\loigiai{
$F_{ht}=mv^2/r=0,5\cdot20^2/5$. Suy ra kết quả $40\ \mathrm{N}$. Đáp án trả lời ngắn không ghi đơn vị.
}
\end{ex}

% ===== Câu 157 =====
% ID: AIQ_L10C6_FULL_0319
% Mức: TH
\begin{bt}
% ID: AIQ_L10C6_FULL_0319
% Mức: TH
Trình bày cách giải: Một vật khối lượng $0,6$ kg quay tròn đều bán kính $2$ m với tốc độ $21$ m/s. Tính lực hướng tâm cần thiết.
\loigiai{
$F_{ht}=mv^2/r=0,6\cdot21^2/2$. Suy ra kết quả $132,3\ \mathrm{N}$.
}
\end{bt}

% ===== Câu 158 =====
% ID: AIQ_L10C6_FULL_0320
% Mức: TH
\begin{bt}
% ID: AIQ_L10C6_FULL_0320
% Mức: TH
Trình bày cách giải: Một vật khối lượng $0,7$ kg quay tròn đều bán kính $3$ m với tốc độ $22$ m/s. Tính lực hướng tâm cần thiết.
\loigiai{
$F_{ht}=mv^2/r=0,7\cdot22^2/3$. Suy ra kết quả $112,93\ \mathrm{N}$.
}
\end{bt}

% ===== Câu 159 =====
% ID: AIQ_L10C6_FULL_0321
% Mức: VD
\begin{bt}
% ID: AIQ_L10C6_FULL_0321
% Mức: VD
Trình bày cách giải: Một vật khối lượng $0,8$ kg quay tròn đều bán kính $4$ m với tốc độ $23$ m/s. Tính lực hướng tâm cần thiết.
\loigiai{
$F_{ht}=mv^2/r=0,8\cdot23^2/4$. Suy ra kết quả $105,8\ \mathrm{N}$.
}
\end{bt}

% ===== Câu 160 =====
% ID: AIQ_L10C6_FULL_0322
% Mức: VD
\begin{bt}
% ID: AIQ_L10C6_FULL_0322
% Mức: VD
Trình bày cách giải: Một vật khối lượng $0,9$ kg quay tròn đều bán kính $5$ m với tốc độ $24$ m/s. Tính lực hướng tâm cần thiết.
\loigiai{
$F_{ht}=mv^2/r=0,9\cdot24^2/5$. Suy ra kết quả $103,68\ \mathrm{N}$.
}
\end{bt}

% ===== Câu 161 =====
% ID: AIQ_L10C6_FULL_0323
% Mức: VDC
\begin{bt}
% ID: AIQ_L10C6_FULL_0323
% Mức: VDC
Trình bày cách giải: Một vật khối lượng $0,5$ kg quay tròn đều bán kính $2$ m với tốc độ $25$ m/s. Tính lực hướng tâm cần thiết.
\loigiai{
$F_{ht}=mv^2/r=0,5\cdot25^2/2$. Suy ra kết quả $156,25\ \mathrm{N}$.
}
\end{bt}

% ===== Câu 162 =====
% ID: AIQ_L10C6_FULL_0324
% Mức: VDC
\begin{bt}
% ID: AIQ_L10C6_FULL_0324
% Mức: VDC
Trình bày cách giải: Một vật khối lượng $0,6$ kg quay tròn đều bán kính $3$ m với tốc độ $26$ m/s. Tính lực hướng tâm cần thiết.
\loigiai{
$F_{ht}=mv^2/r=0,6\cdot26^2/3$. Suy ra kết quả $135,2\ \mathrm{N}$.
}
\end{bt}
